% !TeX program = xelatex
% 数学 LaTeX 排版 Skill v4.0.0；版式保持 v3 金标准. XeLaTeX 编译至少两次. 不分发字体.
% WZ-LOCKED-CLASS-BEGIN
\begin{filecontents*}[overwrite]{elegantbook-original-adapter.cls}
\NeedsTeXFormat{LaTeX2e}
\ProvidesClass{elegantbook-original-adapter}[2026/09/23 v3.0 fixed mathematical book environments]
\LoadClass[10pt,letterpaper,oneside]{book}
\RequirePackage[UTF8,scheme=plain,fontset=fandol]{ctex}
\RequirePackage{fontspec}
\setmainfont{cmunrm.otf}[BoldFont=cmunbx.otf,ItalicFont=cmunti.otf,BoldItalicFont=cmunbi.otf]
\setCJKfamilyfont{sourceheader}[ItalicFont=FandolKai-Regular.otf]{FandolKai-Regular.otf}
\RequirePackage{amsmath,amssymb,mathrsfs}
\RequirePackage{amsthm,etoolbox}
\RequirePackage{xcolor,geometry,setspace,fancyhdr,enumitem,needspace,xurl}
\definecolor{structurecolor}{HTML}{880E4F}
\definecolor{main}{HTML}{9C27B0}
\geometry{paperwidth=612bp,paperheight=792bp,left=20mm,right=20mm,top=95.04bp,bottom=20mm,headheight=12pt,headsep=25pt,footskip=30pt}
\setstretch{1.6}
\setlength{\parindent}{2em}
\setlength{\parskip}{0pt}
\setlist{nosep,leftmargin=2em}
\AtBeginDocument{\setlength{\abovedisplayskip}{4pt plus 1pt minus 1pt}\setlength{\belowdisplayskip}{4pt plus 1pt minus 1pt}\setlength{\abovedisplayshortskip}{2pt}\setlength{\belowdisplayshortskip}{3pt}}
\fancyhf{}
\fancyhead[L]{{\itshape\CJKfamily{sourceheader}\rightmark}}
\fancyhead[R]{{\itshape\CJKfamily{sourceheader}\leftmark}}
\fancyfoot[C]{\thepage}
\renewcommand{\headrulewidth}{0.4pt}
\renewcommand{\headrule}{{\color{structurecolor}\hrule height\headrulewidth width\headwidth\vskip-\headrulewidth}}
\pagestyle{fancy}
\fancypagestyle{cover}{\fancyhf{}\fancyfoot[C]{\thepage}\renewcommand{\headrulewidth}{0pt}\renewcommand{\headrule}{}}
\RequirePackage{hyperref}
\hypersetup{unicode=true,colorlinks=true,linkcolor=structurecolor,urlcolor=structurecolor,citecolor=structurecolor,linktoc=section,bookmarksnumbered=true,bookmarksopen=true,pdfborder={0 0 0},pdfstartview=FitH}
\setcounter{tocdepth}{1}
\renewcommand*\l@section{\@dottedtocline{1}{1.5em}{2.4em}}
\renewcommand{\thesection}{\thechapter.\arabic{section}}
\newcommand{\originalchapter}[1]{\par\addvspace{14pt}\Needspace{5\baselineskip}\refstepcounter{chapter}\setcounter{section}{0}\addcontentsline{toc}{chapter}{\protect\numberline{\thechapter}#1}\markboth{\thechapter\quad #1}{}\begingroup\centering\color{structurecolor}\bfseries\fontsize{14.4}{20}\selectfont\thechapter\quad #1\par\endgroup\nobreak\vspace{8pt}}
\newcommand{\originalsection}[1]{\par\addvspace{9pt}\Needspace{4\baselineskip}\refstepcounter{section}\addcontentsline{toc}{section}{\protect\hspace*{1.5em}\protect\numberline{\protect\textcolor{structurecolor}{\thesection}}#1}\markright{\thesection\quad #1}\noindent{\fontsize{12}{17}\selectfont\color{structurecolor}\thesection\quad}{\bfseries\fontsize{12}{17}\selectfont #1}\par\nobreak\vspace{4pt}}

% Semantic environments are registered once here, not re-designed in manuscripts.
\newtheoremstyle{wzstatement}{6pt}{5pt}
  {\normalfont\kaishu}{0pt}{\normalfont\bfseries\color{structurecolor}}
  {}{.7em}{\thmname{#1}\thmnumber{ #2}\thmnote{\normalfont\ (#3)}}
\newtheoremstyle{wzdefinition}{6pt}{5pt}
  {\normalfont\songti}{0pt}{\normalfont\bfseries\color{structurecolor}}
  {}{.7em}{\thmname{#1}\thmnumber{ #2}\thmnote{\normalfont\ (#3)}}
\newtheoremstyle{wzproblem}{7pt}{5pt}
  {\normalfont\songti}{2em}{\normalfont\bfseries\color{main}}
  {}{.7em}{\thmname{#1}\thmnumber{ #2}}
\newtheoremstyle{wzremark}{4pt}{4pt}
  {\normalfont\songti}{0pt}{\normalfont\bfseries}
  {}{.7em}{\thmname{#1}}
\theoremstyle{wzstatement}
\newtheorem{theorem}{定理}[chapter]
\newtheorem{lemma}[theorem]{引理}
\newtheorem{proposition}[theorem]{命题}
\newtheorem{corollary}[theorem]{推论}
\theoremstyle{wzdefinition}
\newtheorem{definition}[theorem]{定义}
\theoremstyle{wzproblem}
\newtheorem{example}{例题}[chapter]
\newtheorem{exercise}{习题}[chapter]
\theoremstyle{wzremark}
\newtheorem*{remark}{注}
\renewcommand{\proofname}{证明}
% Retain the native amsthm QED stack; only adjust the fixed typography.
\expandafter\patchcmd\csname\string\proof\endcsname{\normalfont \topsep6\p@\@plus6\p@\relax}
  {\normalfont\songti \topsep5\p@\@plus1\p@\relax}
  {}{\ClassError{elegantbook-original-adapter}{Cannot patch proof spacing}{Check the amsthm version.}}
\expandafter\patchcmd\csname\string\proof\endcsname{\itshape}{\normalfont\bfseries}
  {}{\ClassError{elegantbook-original-adapter}{Cannot patch proof heading}{Check the amsthm version.}}
\expandafter\patchcmd\csname\string\proof\endcsname{\ignorespaces}
  {\setlength{\parindent}{2em}\ignorespaces}
  {}{\ClassError{elegantbook-original-adapter}{Cannot patch proof paragraphs}{Check the amsthm version.}}
\AtBeginEnvironment{proof}{\par\Needspace{3\baselineskip}}
\renewcommand{\qedsymbol}{\ensuremath{\square}}
\newenvironment{solution}{\begin{proof}[解答]}{\end{proof}}
\newcommand{\wzcheckchapter}{\ifnum\value{chapter}<1
  \ClassError{elegantbook-original-adapter}{Numbered mathematics before chapter 1}{Move it after the first originalchapter.}\fi}

\newcommand{\wzregisterguard}[1]{\AtBeginEnvironment{#1}{\wzcheckchapter\par\Needspace{3\baselineskip}}}
\forcsvlist{\wzregisterguard}{theorem,lemma,proposition,corollary,definition,example,exercise}
% Front matter and bibliography are not numbered mathematical chapters.
\newcommand{\originalpreface}{\par\addvspace{10pt}\Needspace{5\baselineskip}
  \noindent{\bfseries\fontsize{12}{17}\selectfont 前言}\par\nobreak\vspace{4pt}}
\newcommand{\originalcontents}{\par\addvspace{14pt}
  \begin{center}{\color{structurecolor}\bfseries\fontsize{14.4}{20}\selectfont 目录}\end{center}
  \vspace{-10pt}\@starttoc{toc}}
\renewcommand{\bibname}{参考文献}
\renewenvironment{thebibliography}[1]{
  \par\addvspace{12pt}\Needspace{3\baselineskip}\phantomsection
  \addcontentsline{toc}{chapter}{\bibname}
  \markboth{\bibname}{}
  \noindent{\color{structurecolor}\bfseries\fontsize{14.4}{20}\selectfont\bibname}\par\nobreak\vspace{6pt}
  \list{\@biblabel{\@arabic\c@enumiv}}
    {\settowidth\labelwidth{\@biblabel{#1}}\leftmargin\labelwidth
     \advance\leftmargin\labelsep\usecounter{enumiv}
     \let\p@enumiv\@empty\renewcommand\theenumiv{\@arabic\c@enumiv}
     \setlength{\itemsep}{3pt}\setlength{\parsep}{0pt}}
  \sloppy\clubpenalty4000\widowpenalty4000\sfcode`\.=1000\relax
}{\def\@noitemerr{\@latex@warning{Empty bibliography}}\endlist}

\numberwithin{equation}{chapter}
\renewcommand{\theequation}{\thechapter.\arabic{equation}}
\allowdisplaybreaks[2]
\emergencystretch=2em
\clubpenalty=6000
\widowpenalty=6000
\raggedbottom
\endinput
\end{filecontents*}
% WZ-LOCKED-CLASS-END
\documentclass{elegantbook-original-adapter}
% ===== 元数据内容槽 =====
\newcommand{\BookTitle}{偏微分方程}
\newcommand{\BookAuthor}{Jared Speck 原讲义; 中文重构与增补}
\newcommand{\BookDate}{MIT 18.152, Fall 2011}
\hypersetup{pdftitle={偏微分方程: MIT 18.152 中文重构本},pdfauthor={Jared Speck; 中文整理}}
\usepackage{tikz,booktabs,longtable}
\usetikzlibrary{arrows.meta,calc}
\newcommand{\R}{\mathbb R}
\newcommand{\C}{\mathbb C}
\newcommand{\N}{\mathbb N}
\newcommand{\dd}{\,\mathrm d}
\newcommand{\ee}{\mathrm e}
\newcommand{\ii}{\mathrm i}
\newcommand{\eps}{\varepsilon}
\newcommand{\Sph}{\mathbb S}
\newcommand{\Sch}{\mathcal S}
\newcommand{\supp}{\operatorname{supp}}
\newcommand{\diver}{\operatorname{div}}
\newcommand{\sourcelecture}[1]{\par\noindent 对应原课程第 #1 讲.\par}
\begin{document}
\frontmatter
\begin{center}
{\color{structurecolor}\bfseries\fontsize{18}{25}\selectfont\BookTitle\par}
{\bfseries MIT 18.152 中文重构本\par}
\BookAuthor\par
\BookDate\par
\end{center}
\originalpreface
热量的扩散、静电势的平衡和声波的传播都可以写成偏微分方程, 但相同的求导符号并不意味着相同的解法. 本书围绕初值、边界值、解的意义以及能量估计, 建立热方程、Laplace/Poisson 方程和波方程的完整入门理论, 再以 Fourier 变换连接扩散与色散, 进入 Schrödinger 方程、变分场论和输运方程.

读者需熟悉一元和多元微积分、线性代数、常微分方程以及基础实分析. 散度定理、控制收敛、Fubini 定理和 $L^2$ 完备性作为先修工具使用, 其应用条件在具体证明中说明. 分布、弱解、Fourier 归一化及能量空间在首次使用处定义. 不要求先学 Sobolev 空间的一般理论.

依据 Jared Speck 的 \href{https://ocw.mit.edu/courses/18-152-introduction-to-partial-differential-equations-fall-2011/pages/lecture-notes/}{MIT OpenCourseWare 18.152 Introduction to Partial Differential Equations, Fall 2011} 全部18份公开课堂讲义, 共24课时、136个实际 PDF 页面, 含18页末尾 OCW 来源页. 本书按依赖关系重排, 合并反复出现的记号、守恒律和复习内容, 不嵌入英文页面. 定义、定理、证明、例题和解答均为简体中文. 公式采用统一的符号约定, 补齐原讲义留作作业的核心论证, 并在末章区分原课程内容、编者补充与未覆盖专题.

官方英文原件仅存于仓库的 \texttt{sources/partial-differential-equations/}, 不收入中文源码 ZIP. 原作者 Jared Speck, 原机构 Massachusetts Institute of Technology. 中文翻译、章节重组、条件修订、增补证明、自绘图与训练题构成本整理本的改动, 不表示 MIT 或原作者认可本项目. 官方资源页及各 PDF 的署名、许可提示、页数和校验和逐项登记于来源清单; 未发现单独的第三方复制保留声明. 数学书目引用不被解释为所引书籍可以再分发.

MIT OCW 当前默认许可为 \href{https://creativecommons.org/licenses/by-nc-sa/4.0/}{CC BY-NC-SA 4.0 International}, 本中文改编沿用相同许可. 原件的末页指向 OCW 使用条款, 单独权利声明优先于默认许可. 条款见 \href{https://ocw.mit.edu/pages/privacy-and-terms-of-use/}{MIT OCW 使用条款}. 本册同时收入全部公开完整作业题、Bonus、期中和期末试题的中文题面与解答. 期中答案依据官方解答重构并修订, 其余答案由 AI 独立编写并交叉审读. 仅有题号的11条商业教材引用及1条阅读任务登记在缺项表, 不猜造题面, 不收入商业教材. 版本日期: 2026-10-08.
\clearpage
\phantomsection
\addcontentsline{toc}{chapter}{原课程说明与进度}
\markboth{原课程说明与进度}{}
\begin{center}{\color{structurecolor}\bfseries\fontsize{14.4}{20}\selectfont 原课程说明与进度}\end{center}
本节译自官方 \href{https://ocw.mit.edu/courses/18-152-introduction-to-partial-differential-equations-fall-2011/pages/syllabus/}{教学大纲}与\href{https://ocw.mit.edu/courses/18-152-introduction-to-partial-differential-equations-fall-2011/pages/calendar/}{课程日历}, 不把本整理本的章节顺序解释为原课日历.

课程为 MIT 数学系本科课18.152, 2011年秋季, 教师 Jared Speck. 每周两次课, 每次1.5小时. 先修18.100分析, 以及18.06线性代数、18.700线性代数或18.701代数 I 三者之一. 原课围绕热扩散、静电势和波传播建立方法, 再学习 Fourier 变换、Schrödinger 方程与 Lagrange 场论. 大纲提及 Maxwell 方程等视时间安排的专题; 公开讲义没有建立 Maxwell 系统理论.

指定书为 Sandro Salsa, \textit{Partial Differential Equations in Action: From Modelling to Theory}, Springer, 2010, ISBN9788847007512. 这是商业书目, 不是本项目可再分发的公开原件. 几乎每周一次作业, 通常周四发布、下一周四交, 不接受迟交, 最低一次作业成绩剔除. 期中90分钟, 期末综合全课且稍侧重后半段; 大纲称均闭卷、不得用笔记; 实际期末试卷允许一张正反面手写笔记, 见期末章说明, 以原卷为准. 总评作业30\%, 期中30\%, 期末40\%. 原日历仅列课序, 未列逐课公历日期; 下表不推测日期.

\begin{longtable}{p{.12\textwidth}p{.58\textwidth}p{.21\textwidth}}
\toprule 课序&官方主题的中文译名&提交/考试\\\midrule\endhead
L1&偏微分方程导论&\\
L2&热方程导论&\\
L3&热方程唯一性&作业1\\
L4&热方程弱最大原理及基本解导论&\\
L5&热基本解与全空间 Cauchy 问题&作业2\\
L6&Laplace 与 Poisson 方程&\\
L7&Poisson 基本解&作业3\\
L8&Poisson 的 Green 函数&\\
L9&Poisson 公式、Harnack 不等式、Liouville 定理&作业4\\
L10&波方程导论&作业5\\
L11&波方程球面平均法&\\
L12&Kirchhoff 公式与 Minkowski 几何&作业6\\
L13&波方程几何能量估计&\\
E1&期中考试&90分钟\\
L14&波方程几何能量估计续&\\
L15&二阶方程分类&作业7\\
L16&Fourier 变换导论&\\
L17&Fourier 变换续&作业8\\
L18&Fourier 反演与 Plancherel 定理&\\
L19&Schrödinger 方程导论&作业9\\
L20&Schrödinger 方程续&\\
L21&Lagrange 场论导论&选做 Bonus\\
L22&Lagrange 场论续&作业10\\
L23&Lagrange 场论续&\\
L24&输运与 Burgers 方程&作业11\\
E2&期末考试&综合全课\\\bottomrule
\end{longtable}

\clearpage
\markboth{目录}{}
\originalcontents
\clearpage
\mainmatter
\originalchapter{方程、数据与解的意义}\label{ch:basic}
\sourcelecture{1}
\originalsection{偏微分算子与线性结构}
设 $\Omega\subset\R^n$ 是开且连通的区域. 多重指标 $\alpha=(\alpha_1,\ldots,\alpha_n)\in\N^n$ 的阶为 $|\alpha|=\sum_j\alpha_j$, $\partial^\alpha=\partial_1^{\alpha_1}\cdots\partial_n^{\alpha_n}$. 对含时间的函数写 $u(t,x)$, 空间梯度为 $\nabla u=(u_{x_1},\ldots,u_{x_n})$, Laplace 算子为 $\Delta u=\sum_j u_{x_jx_j}$. 时间不计入 $\Delta$.
\begin{definition}
含未知函数及其偏导数的关系 $F(x,u,\partial u,\ldots,\partial^m u)=0$ 称为偏微分方程. 实际出现的最高导数阶数称为阶. 若 $L=\sum_{|\alpha|\le m}a_\alpha(x)\partial^\alpha$, 则 $Lu=f$ 为线性方程; $f=0$ 时称为齐次. 系数依赖于未知函数时一般不再线性.
\end{definition}
热方程 $u_t-D\Delta u=0$ 是时间一阶、空间二阶, 按最高导数阶数是二阶方程. 波方程 $u_{tt}-c^2\Delta u=0$ 是二阶方程. $u_t+uu_x=0$ 是一阶非线性方程, 非线性来自 $u$ 与 $u_x$ 的乘积. 齐次指右端为零, 与两端边界条件是否同一种类型是不同概念.
\begin{proposition}\label{prop:superposition}
线性算子 $L$ 的齐次解组成向量空间. 若 $Lu_p=f$, 则全部非齐次解恰为 $u_p+v$, 其中 $Lv=0$.
\end{proposition}
\begin{proof}
由 $L(au+bv)=aLu+bLv$ 得齐次解对线性组合封闭. 若 $Lu=f$, 则 $L(u-u_p)=0$; 反过来 $L(u_p+v)=f+0=f$. 这同时证明两类解的对应. 若还规定边界数据, 相加时数据也必须相加, 不能把任意齐次方程解加入而保持原边界值.
\end{proof}
\begin{example}
求 $a u_x+b u_y=0$ 的解, 其中 $(a,b)\ne(0,0)$, 并分析数据 $u(x,0)=g(x)$.
\end{example}
\begin{solution}
沿 $(x,y)=(x_0+as,y_0+bs)$ 有 $\frac{\mathrm d}{\mathrm ds}u=a u_x+b u_y=0$. 取横向坐标 $r=bx-ay$, 则每条直线上 $r$ 不变, 在全平面全部 $C^1$ 解是 $u=F(r)$. 若 $b\ne0$, 数据给 $F(bx)=g(x)$, 故 $u(x,y)=g(x-ay/b)$. 若 $b=0$, 数据线本身为特征线, 数据必须是常数, 而不同水平线上的值仍可任意选择. 数据并非越多越好, 要与方程的传播方向相容.
\end{solution}
\originalsection{经典解、分布与弱解}
\begin{definition}
经典解是具有方程所需连续导数并逐点满足方程的函数. 对热方程采用 $C^{1,2}$: 函数、一次时间导数和至多二次空间导数连续. 边界与初值是否逐点成立, 需另行规定. $L^p(\Omega)$ 由可测函数的几乎处处等价类组成, 范数为 $\|u\|_p=(\int_\Omega |u|^p\dd x)^{1/p}$, $p<\infty$, 而 $\|u\|_\infty$ 是本质上确界.
\end{definition}
一次积分分部常比逐点二次导数要求更弱. 记 $C_c^\infty(\Omega)$ 为具有紧支撑的光滑测试函数. 分布是其上的连续线性泛函, 局部可积函数 $u$ 定义 $\langle u,\phi\rangle=\int u\phi$. 分布导数定义为 $\langle\partial_j T,\phi\rangle=-\langle T,\partial_j\phi\rangle$. 这里连续指测试函数支撑位于同一紧集且所有阶导数一致趋零时, 泛函值也趋零. Dirac 分布满足 $\langle\delta_a,\phi\rangle=\phi(a)$, 不是普通函数.
\begin{definition}
局部可积的 $u$ 是 $-\Delta u=f$ 的分布解, 若对每个 $\phi\in C_c^\infty(\Omega)$ 有 $\int u(-\Delta\phi)=\int f\phi$. 若 $u$ 具有平方可积的弱一阶导数, 则等价的弱形式是 $\int\nabla u\cdot\nabla\phi=\int f\phi$. 初边值问题的弱解还必须指定初始迹和边界条件, 仅写内点测试等式并不完整.
\end{definition}
\begin{example}
计算 $u(x)=|x|$ 的分布二阶导数.
\end{example}
\begin{solution}
把积分在零处分开, 对紧支撑 $\phi$ 积分分部两次:
\[
\int_\R |x|\phi''(x)\dd x
=\int_{-\infty}^0(-x)\phi''\dd x+\int_0^\infty x\phi''\dd x=2\phi(0).
\]
故 $u''=2\delta_0$. 零点以外的经典导数为零, 却不能描述折角处的集中源. 这正是引入弱导数的作用.
\end{solution}
\originalsection{适定性与积分恒等式}
\begin{definition}
在指定数据空间和解空间中, 一个问题称为适定, 若每个允许数据都有解, 解唯一, 且解连续依赖数据. 连续依赖必须带有范数及时间区间; 仅有显式公式不构成适定性证明.
\end{definition}
热方程需给一个初值, 波方程需给位移和速度两个初值. 在有界区域还须选边界条件. Dirichlet 给 $u$, Neumann 给外法向导数 $\partial_\nu u=\nabla u\cdot\nu$, Robin 给 $\partial_\nu u+\alpha u$. 一维区间左端外法向是 $-1$, 右端是 $1$, 因而左端 $\partial_\nu u=-u_x$, 右端 $\partial_\nu u=u_x$.

对 $C^1$ 边界的有界区域及闭包上足够光滑的函数, 散度定理给出
\begin{align}
\int_\Omega v\Delta u\dd x
&=\int_{\partial\Omega}v\partial_\nu u\dd S-\int_\Omega\nabla v\cdot\nabla u\dd x,\label{eq:greenone}\\
\int_\Omega(v\Delta u-u\Delta v)\dd x
&=\int_{\partial\Omega}(v\partial_\nu u-u\partial_\nu v)\dd S.\label{eq:greentwo}
\end{align}
第一式对向量场 $v\nabla u$ 应用散度定理; 再交换 $u,v$ 相减即得第二式. 后文反复用这两个等式: 能量法令 $v=u$, Green 函数法令 $v$ 为带点源的辅助函数. 对奇异核不能直接套用, 必须先挖去小球再取极限.

积分号下求导采用以下先修判据: 参数邻域内被积函数可微, 导数被同一个可积函数控制, 则可交换求导与积分. 证明由差商、微积分基本定理和控制收敛给出. 后文热核的导数是多项式乘 Gaussian, 正时间紧区间内提供所需控制函数.
\begin{exercise}
验证 $u_{xx}+u_{yy}=0$ 的所有仿射函数解, 并说明只给 $u(0,0)=0$ 为什么不能确定解.
\end{exercise}
\begin{solution}
$u=ax+by+d$ 的二阶导数为零. 原点条件仅令 $d=0$, 仍留下两个自由参数. 事实上还有无限多其他调和函数; 单点值不是二维区域边界值问题所需的数据.
\end{solution}

\originalchapter{热方程的边界问题与能量}\label{ch:heatboundary}
\sourcelecture{2--4}
\originalsection{守恒律与边界条件}
温度 $u$, 密度 $\rho>0$, 比热 $c_v>0$ 和热导率 $\kappa>0$ 均在均匀各向同性模型中使用. 任意子区域 $V$ 的热能为 $\int_V\rho c_vu$. 若单位质量热源率为 $R$, 热流为 $q=-\kappa\nabla u$, 则
\[
\frac{\mathrm d}{\mathrm dt}\int_V\rho c_vu\dd x
=-\int_{\partial V}q\cdot\nu\dd S+\int_V\rho R\dd x
=\int_V(\kappa\Delta u+\rho R)\dd x.
\]
连续被积函数在任意小区域上积分为零意味着自身为零, 故 $u_t-D\Delta u=f$, 其中 $D=\kappa/(\rho c_v)$, $f=R/c_v$. 这是模型推导, Fourier 热传导定律及常比热假设有其物理适用范围.

令 $Q_T=(0,T]\times\Omega$. 抛物边界为 $\partial_pQ_T=(\{0\}\times\overline\Omega)\cup([0,T]\times\partial\Omega)$, 不包括最后时刻的内部. Dirichlet 数据 $u(0,x)=g(x)$, $u(t,\xi)=h(t,\xi)$ 若要求闭柱连续, 则角点须有 $g|_{\partial\Omega}=h(0,\cdot)$. 更高正则性还需更高相容条件; 例如 $g\in C^2$ 时可能要 $D\Delta g+f(0)=h_t(0)$ 于边界.
\originalsection{能量耗散、唯一性与稳定性}
\begin{theorem}\label{thm:heatenergy}
设 $\Omega$ 有界、边界 $C^1$, $u$ 为闭柱上 $C^{1,2}$ 的齐次热方程解. 边界为齐次 Dirichlet、齐次 Neumann 或 $\partial_\nu u+\alpha u=0$、$\alpha\ge0$ 的 Robin 条件时,
\[
\frac12\frac{\mathrm d}{\mathrm dt}\|u(t)\|_2^2
=-D\|\nabla u(t)\|_2^2-D\int_{\partial\Omega}\alpha|u|^2\dd S\le0,
\]
其中 Dirichlet/Neumann 情形最后一项取零. 分段混合条件也成立, 每段采用相应项. 同一初值和边界数据的解唯一.
\end{theorem}
\begin{proof}
将 $u_t=D\Delta u$ 乘 $u$ 积分, 用\eqref{eq:greenone}. Dirichlet 段因 $u=0$ 消去边界项, Neumann 段因 $\partial_\nu u=0$ 消去, Robin 段为 $-\alpha u^2$. 取两解之差 $w$, 它有零初值及齐次边界, 故 $0\le\|w(t)\|_2^2\le\|w(0)\|_2^2=0$. 连续性使几乎处处为零强化为处处为零.
\end{proof}
Robin 的符号不能任意省略. 负 $\alpha$ 会使上述耗散证明失效, 不等于立即失去唯一性; 在合适条件下仍可用边界估计与 Gronwall 不等式处理, 本册使用 $\alpha\ge0$ 的版本.
\begin{proposition}\label{prop:heatl2stable}
若两解之差满足 $w_t-D\Delta w=F$ 与相同类型的齐次耗散边界, 则
\[
\|w(t)\|_2\le\|w(0)\|_2+\int_0^t\|F(s)\|_2\dd s.
\]
\end{proposition}
\begin{proof}
能量等式与 Cauchy--Schwarz 给 $\frac12(\|w\|_2^2)'\le\|w\|_2\|F\|_2$. 为避免在 $\|w\|_2=0$ 除零, 设 $y_\eps=(\|w\|_2^2+\eps^2)^{1/2}$, 得 $y_\eps'\le\|F\|_2$. 积分后令 $\eps\downarrow0$ 即得结果.
\end{proof}
\begin{example}
绝热杆 $u_t=Du_{xx}$, $u_x(t,0)=u_x(t,L)=0$ 的平均温度如何变化?
\end{example}
\begin{solution}
$\frac{\mathrm d}{\mathrm dt}\int_0^Lu\dd x=D[u_x]_0^L=0$, 所以平均 $\overline u=L^{-1}\int_0^Lg$ 守恒. 能量耗散的是偏离平均的部分, 并不强迫非零平均趋于零. 若给外法向导数 $h_0,h_L$, 则总量变化为 $D(h_0+h_L)$; 左端符号已包括在外法向定义里.
\end{solution}
\originalsection{弱最大值原理}
\begin{theorem}\label{thm:heatmax}
$\Omega$ 有界. 设 $w\in C(\overline Q_T)\cap C^{1,2}((0,T)\times\Omega)$, 且 $w_t-D\Delta w\le0$ 于内部. 则 $\max_{\overline Q_T}w\le\max_{\partial_pQ_T}w$.
\end{theorem}
\begin{proof}
先在 $[0,\tau]\times\overline\Omega$, $\tau<T$, 上考察 $v=w-\eps t$. 若最大点有 $t_0>0$, $x_0\in\Omega$, 空间 Hessian 半负定, 故 $\Delta v\le0$; 时间内点 $v_t=0$, 上顶点以左导数得 $v_t\ge0$. 因而 $v_t-D\Delta v\ge0$, 与 $w_t-D\Delta w-\eps<0$ 矛盾. 所以 $v$ 的最大值在抛物边界, 从而 $w\le M+\eps\tau$, 其中 $M=\max_{\partial_pQ_T}w$. 令 $\eps\downarrow0$, 再令 $\tau\uparrow T$, 连续性给闭柱上的结论.
\end{proof}
\begin{corollary}\label{cor:heatcompare}
若 $u_t-D\Delta u=f$, $v_t-D\Delta v=g$, 且抛物边界上 $u\le v$, 内部 $f\le g$, 则 $u\le v$. 此外,
\[
\|u-v\|_{C(\overline Q_T)}\le\|u-v\|_{C(\partial_pQ_T)}+T\|f-g\|_{C(\overline Q_T)}.
\]
\end{corollary}
\begin{proof}
对 $u-v$ 应用最大值原理得比较结论. 令 $A=\|u-v\|_{\partial_pQ_T}$, $B=\|f-g\|_\infty$. 对 $u-v-A-Bt$ 与 $v-u-A-Bt$ 分别应用原理, 得 $|u-v|\le A+Bt$.
\end{proof}
取 $v=0$ 可知非负初边值与非负源产生非负解. 这给出顺序结构及一致范数稳定性. 定理要求空间有界; 在全空间中最大值可能不取得, 后章用增长控制处理无穷远.
\begin{exercise}
若 $u_t-D\Delta u+qu=0$, $q\ge0$ 为连续函数, 初边值满足 $0\le u\le M$, $M\ge0$, 证明 $0\le u\le M$.
\end{exercise}
\begin{solution}
对该算子作严格扰动最大值证明, 在正的内部最大点有 $u_t-D\Delta u+qu\ge0$, 因 $qu\ge0$, 与严格负右端矛盾. 比较 $u$ 与零, 得下界; 常数 $M$ 满足 $M_t-D\Delta M+qM=qM\ge0$, 是上解, 故得上界. 可用 $\eps\ee^{\lambda t}$ 作严格扰动, $\lambda>0$, 在有限柱上确保不等式严格.
\end{solution}

\originalchapter{分离变量与 Fourier 级数}\label{ch:series}
\sourcelecture{2--3; Fourier 分析所需的论证为编者补充}
\originalsection{Dirichlet 特征函数}
在杆 $(0,L)$ 上令 $u(t,x)=T(t)X(x)$, 代入 $u_t=Du_{xx}$, 非零因子处有 $T'/(DT)=X''/X=-\lambda$. 因而
\[
-X''=\lambda X,\quad X(0)=X(L)=0,\qquad T'= -D\lambda T.
\]
乘 $X$ 积分得 $\lambda\int X^2=\int (X')^2$, 非零解要求 $\lambda>0$. 解 ODE 后边界条件给 $\lambda_k=(k\pi/L)^2$, $X_k(x)=\sin(k\pi x/L)$, $k\ge1$. 零特征值与负特征值只有零解. 正交关系是
\[
\int_0^L\sin(k\pi x/L)\sin(j\pi x/L)\dd x=\frac L2\delta_{kj}.
\]
由积化和差公式积分即可验证. 分离变量仅找到了单个模态; 要表示任意数据, 还需证明这些模态在 $L^2$ 中完备.
\originalsection{完备性与初始迹}
以下给出所需的完备性论证, 避免把有限叠加原理直接用于未经检验的无穷和. 在周期 $2\pi$ 上, 定义 Fejér 核
\[
F_N(\theta)=\frac1{N+1}\left|\sum_{j=0}^N\ee^{\ii j\theta}\right|^2
=\sum_{|k|\le N}\left(1-\frac{|k|}{N+1}\right)\ee^{\ii k\theta}.
\]
它非负, $\int_{-\pi}^{\pi}F_N\dd\theta=2\pi$, 且 $|\theta|\ge\delta>0$ 时由几何级数得 $F_N\le[(N+1)\sin^2(\delta/2)]^{-1}$. 所以质量集中在零点附近.
\begin{lemma}\label{lem:fejer}
对周期 $L^2$ 函数 $f$, $\sigma_N f=(2\pi)^{-1}F_N*f$ 在 $L^2$ 中趋于 $f$. 对连续周期函数, 收敛是一致的.
\end{lemma}
\begin{proof}
周期平移在 $L^2$ 中连续: 阶梯函数可直接算出平移造成的误差, 再用阶梯函数在 $L^2$ 中稠密及平移保范数, 得 $\|f(\cdot-\theta)-f\|_2\to0$. 由积分 Minkowski 不等式,
\[
\|\sigma_N f-f\|_2\le\frac1{2\pi}\int_{-\pi}^{\pi}F_N(\theta)\|f(\cdot-\theta)-f\|_2\dd\theta.
\]
在 $|\theta|<\delta$ 用平移连续性使第二因子小于 $\eps$, 其余部分用 $2\|f\|_2$ 及核在远处的质量趋零. 故上极限至多 $\eps$. 连续函数用一致连续性和上确界范数重复同一论证.
\end{proof}
\begin{theorem}\label{thm:sinebasis}
$e_k=(2/L)^{1/2}\sin(k\pi x/L)$ 为 $L^2(0,L)$ 的完备正交规范系. 对 $g\in L^2$, 写 $b_k=(2/L)\int_0^Lg(x)\sin(k\pi x/L)\dd x$, 则
\[
g=\sum_{k=1}^\infty b_k\sin(k\pi x/L)\quad\text{于 }L^2,
\qquad \|g\|_2^2=\frac L2\sum_{k=1}^\infty |b_k|^2.
\]
\end{theorem}
\begin{proof}
把 $g$ 奇延拓到 $(-L,L)$ 并周期延拓. 若 $g$ 与所有正弦正交, 奇延拓的余弦及常数系数因奇性为零, 正弦系数也为零. Fejér 平均全为零, 引理迫使延拓函数为零, 所以正交补为零. 有限正交投影的误差与其余子空间正交, 勾股定理给 Bessel 不等式. 系数平方可和使投影在 $L^2$ 中 Cauchy; 完备性给极限, 其误差与全部 $e_k$ 正交, 故为零. 取范数得 Parseval 等式.
\end{proof}
连续函数的普通 Fourier 部分和并不总是一致收敛. 本册使用已证明的 $L^2$ 完备性和 Fejér 一致逼近, 不采用“连续即普通级数一致收敛”的过强说法.
\originalsection{区间热方程的构造}
\begin{theorem}\label{thm:intervalheat}
$g\in L^2(0,L)$ 时,
\[
u(t,x)=\sum_{k\ge1}b_k\ee^{-D\lambda_kt}\sin(k\pi x/L)
\]
在每个 $t\ge\tau>0$ 上可逐项任意求导, 满足热方程及零 Dirichlet 边界, 且 $u(t)\to g$ 于 $L^2$ 当 $t\downarrow0$. 在 $C([0,T];L^2)$、正时间光滑且可作能量积分的此类解中唯一, 并有
\[
\|u(t)\|_2\le\ee^{-D\pi^2t/L^2}\|g\|_2.
\]
\end{theorem}
\begin{proof}
任意时间、空间导数使第 $k$ 项至多多一个 $k^m$ 因子. Cauchy--Schwarz 和 $\sum k^{2m}\ee^{-2D\lambda_k\tau}<\infty$ 保证一致绝对收敛, 故可逐项求导. 每项都满足方程和边界. Parseval 给
\[
\|u(t)-g\|_2^2=\frac L2\sum_k|b_k|^2|\ee^{-D\lambda_kt}-1|^2\longrightarrow0,
\]
用可和控制项 $4|b_k|^2$ 交换极限. 同式给衰减估计. 唯一性在 $[\eps,t]$ 上用能量法得 $\|w(t)\|_2\le\|w(\eps)\|_2$, 令 $\eps\downarrow0$.
\end{proof}
\begin{example}\label{ex:rodx}
求 $L=D=1$, $g(x)=x$, 零 Dirichlet 边界的热方程解, 并解释右下角的不连续.
\end{example}
\begin{solution}
积分分部给 $b_k=2\int_0^1x\sin(k\pi x)\dd x=2(-1)^{k+1}/(k\pi)$. 故
\[
u(t,x)=\sum_{k\ge1}\frac{2(-1)^{k+1}}{k\pi}\ee^{-k^2\pi^2t}\sin(k\pi x).
\]
初值在 $L^2$ 意义成立. 对内部点的点态初始迹, 将 $g(x)=x$ 奇延拓并以周期2延拓, 得有界锯齿函数 $G$. 周期化热核 $K_{\rm per}(t,z)=\sum_{m\in\mathbb Z}K_1(t,z+2m)$ 在一个周期内质量为1, 非负, 远离零模周期的质量趋零. 因而周期卷积在 $G$ 的连续内点趋 $G(x)$. 对周期模态直接积分 Gaussian 得因子 $\ee^{-k^2\pi^2t}$, 由 $L^2$ 完备性, 该卷积正是上面的级数. 这给 $0<x<1$ 的点态初始迹, 不借助未证明的普通 Fourier 级数收敛断言. $t>0$ 时 $u(t,1)=0$, 而 $g(1)=1$, 因而不能在 $(0,1)$ 角点连续. 一个点不影响 $L^2$ 初始迹, 所以这不妨碍定理的存在和唯一性.
\end{solution}
Neumann 情形用 $1,\cos(k\pi x/L)$; 常数模态不衰减, 其系数是平均温度. 非零恒定 Dirichlet 边界可先减去连接端点值的仿射稳态函数. 随时间变化的边界可减去同样的仿射函数 $H(t,x)$, 但新方程的源必须改成 $f-H_t+D H_{xx}$.
\begin{exercise}
解 $u_t=Du_{xx}$, $u(t,0)=A$, $u(t,L)=B$, $u(0,x)=g(x)$, 其中 $g$ 连续且满足角点相容条件.
\end{exercise}
\begin{solution}
令 $H(x)=A+(B-A)x/L$, $v=u-H$, 则 $v$ 满足零边界和初值 $g-H$. 对 $g-H$ 展开正弦级数得到上述解, 再加 $H$. 能量估计给 $\|u(t)-H\|_2\le\ee^{-D\pi^2t/L^2}\|g-H\|_2$.
\end{solution}

\originalchapter{热核与全空间初值问题}\label{ch:heatkernel}
\sourcelecture{4--5}
\originalsection{缩放与 Gaussian 核}
抛物缩放 $u_\lambda(t,x)=\lambda^nu(\lambda^2t,\lambda x)$ 保持方程, 也保持空间积分. 集中质量随时间扩散的尺度应为 $|x|\sim\sqrt{Dt}$. 一维试取 $K(t,x)=(Dt)^{-1/2}V(\zeta)$, $\zeta=x/\sqrt{Dt}$, 得
\[
V''+\tfrac12\zeta V'+\tfrac12V=0,
\quad (V'+\tfrac12\zeta V)'=0.
\]
取偶对称且衰减的解, 在零点有 $V'(0)=0$, 故 $V'+\zeta V/2=0$, 得 $V=C\ee^{-\zeta^2/4}$. Gaussian 积分由二维极坐标求出 $\int_\R\ee^{-s^2}\dd s=\sqrt\pi$. 归一化后, 各方向相乘给
\begin{equation}\label{eq:heatkernel}
K_D(t,x)=(4\pi Dt)^{-n/2}\ee^{-|x|^2/(4Dt)},\qquad t>0.
\end{equation}
直接求导有 $\partial_tK_D=(-n/(2t)+|x|^2/(4Dt^2))K_D$ 及 $\Delta K_D=(-n/(2Dt)+|x|^2/(4D^2t^2))K_D$, 所以 $K_t=D\Delta K$. 换元 $x=\sqrt{4Dt}\,z$ 得 $\int K_D=1$. 注意 $x=0$ 时核趋于无穷, $x\ne0$ 时趋零; 所谓初值 $\delta_0$ 指分布极限.
\begin{lemma}\label{lem:heatapprox}
若 $g$ 有界连续, $K_D(t)*g(x)\to g(x)$ 在每个连续点成立; 若 $g$ 有界且一致连续, 则一致收敛. 对 $1\le p<\infty$, $g\in L^p$ 时收敛于 $L^p$.
\end{lemma}
\begin{proof}
卷积定义为 $(K*g)(x)=\int K(y)g(x-y)\dd y$. 减去 $g(x)\int K$ 后, 在 $|y|<\delta$ 用连续性控制差值, 其余部分用 $2\|g\|_\infty\int_{|y|\ge\delta}K$. 换元后尾积分下限为 $\delta/\sqrt{4Dt}\to\infty$, 故趋零. 一致连续使 $\delta$ 不依赖 $x$. $L^p$ 情形以 $\|g(\cdot-y)-g\|_p$ 替代点态差值, 利用平移连续性和积分 Minkowski 不等式. 平移连续性由紧支撑阶梯函数逼近证明, 与引理\ref{lem:fejer}相同.
\end{proof}
\originalsection{构造、正性与无穷传播}
\begin{theorem}\label{thm:heatrn}
$g$ 有界连续时, $u(t)=K_D(t)*g$ 是全空间热方程的有界经典解, 正时间 $C^\infty$, 连续取得初值. 有界解中唯一, 且 $\|u(t)\|_\infty\le\|g\|_\infty$. $g\in L^p$ 时同公式满足 $\|u(t)\|_p\le\|g\|_p$ 及 $L^p$ 初始迹, $p<\infty$.
\end{theorem}
\begin{proof}
在 $t\in[\tau,T]$, $x$ 位于紧集时, 核的任意导数均被 $C(1+|y|)^m\ee^{-c|y|^2}$ 控制, 可积. 所以可交换求导与积分, 用 $K_t=D\Delta K$ 得方程和光滑性. 初值由上引理. 核非负且积分为1, 给一致范数估计. 积分 Minkowski 与平移保 $L^p$ 范数给 $L^p$ 收缩.

证明有界唯一性, 取零初值解 $w$ 及任意 $\eps>0$. 设 $\Phi(t,x)=|x|^2+2nDt+1$, 则 $\Phi_t-D\Delta\Phi=0$. 在大球 $B_R$ 的侧边, 有界性使 $w-\eps\Phi\le0$; 初时也非正. 最大值原理给 $w\le\eps\Phi$ 于该柱. 对固定点先取足够大 $R$, 再令 $\eps\downarrow0$, 得 $w\le0$. 对 $-w$ 重复即得零.
\end{proof}
$g\ge0$ 且非零于正测度集合时, $K_D>0$ 使 $u(t,x)>0$ 对每个 $t>0,x$ 成立. 所以紧支撑的非负热量立即影响全空间. 对一般带符号数据, 不能仅从核正性推出每点非零. $L^1$ 数据的总量由 Fubini 保持: $\int u=\int g$. 全空间的长期衰减通常为幂次, 如 $\|u(t)\|_\infty\le(4\pi Dt)^{-n/2}\|g\|_1$, 与有界杆的谱隙导致的指数衰减不同.
\begin{example}
一维初值 $g(x)=\ee^{-a x^2}$, $a>0$, 求热方程解并计算总量.
\end{example}
\begin{solution}
配方积分得到
\[
u(t,x)=(1+4aDt)^{-1/2}\exp\left(-\frac{a x^2}{1+4aDt}\right).
\]
宽度增大而峰值减小, 积分始终为 $\sqrt{\pi/a}$. 可直接求导核验 PDE, 也可由卷积构造的唯一性认定.
\end{solution}
\originalsection{增长数据与唯一性类别}
\begin{proposition}\label{prop:growthheat}
若 $g$ 连续且 $|g(x)|\le A\ee^{b|x|^2}$, $b>0$, 则核卷积在 $0<t<(4Db)^{-1}$ 收敛并给经典解. 在 $[0,T']$, $T'<(4Db)^{-1}$, 有
\[
|u(t,x)|\le A(1-4Dbt)^{-n/2}\exp\left(\frac{b|x|^2}{1-4Dbt}\right).
\]
在每个有限时间柱上满足某个统一 Gaussian 增长界的解中, 初值解唯一.
\end{proposition}
\begin{proof}
将 $b|y|^2-|x-y|^2/(4Dt)$ 配方, 二次项系数为负恰当且仅当 $4Dbt<1$. 积分得到所示界. 正时间紧区间内导数仍是多项式乘严格衰减 Gaussian, 可求导. 初始迹在近点用连续性, 尾部以 $\ee^{b|y|^2}$ 和小时间核相乘的衰减控制, 得 $u(t,x)\to g(x)$.

对两个增长解之差, 在固定时间区间取共同界 $|w|\le C\ee^{B|x|^2}$. 选 $a>B$ 和小时间 $0\le t\le\tau<(4Da)^{-1}$. 函数
\[
\Psi(t,x)=(1-4Dat)^{-n/2}\exp\left(\frac{a|x|^2}{1-4Dat}\right)
\]
是热方程解, 由同一配方或直接求导核实, 且 $\Psi\ge\ee^{a|x|^2}$. 对 $w-\eps\Psi$ 在 $B_R$ 用最大值原理; $C\ee^{BR^2}\le\eps\ee^{aR^2}$ 对充分大 $R$ 成立. 先取 $R$, 再令 $\eps\downarrow0$, 得 $w\le0$, 同理 $-w\le0$. 用同样长度的短时间段从已知零截面重新开始, 有限次覆盖整个区间.
\end{proof}
唯一性始终包含解类别. 放弃无穷远增长限制会出现零初值的非平凡热方程解, 本册不建立该反例的一般构造, 也不声称所有光滑全空间解唯一.
\originalsection{Duhamel 原理与半空间反射}
\begin{theorem}\label{thm:duhamelheat}
设 $g$ 有界连续, $f$ 及其空间一、二阶导数在 $[0,T]\times\R^n$ 有界连续. 则
\[
u(t)=K_D(t)*g+\int_0^tK_D(t-s)*f(s)\dd s
\]
是 $u_t-D\Delta u=f$ 的有界解, 正时间 $C^{1,2}$, 初值为 $g$.
\end{theorem}
\begin{proof}
记积分项 $v$. 为处理 $s=t$ 的核奇性, 将空间导数移到 $f$, 得 $\Delta v=\int_0^tK_D(t-s)*\Delta f(s)\dd s$. 积分分部的无穷远边界项因核衰减与 $f,\nabla f$ 有界而消失. 用 $K_D(r)*f(s)\to f(s)$ 及二阶导数的有界控制, Leibniz 法则给
\[
v_t=f(t)+\int_0^tD K_D(t-s)*\Delta f(s)\dd s=f+D\Delta v.
\]
$\|v(t)\|_\infty\le t\|f\|_\infty$ 给零初值. 连续性由近似单位和控制收敛保证. 两解之差的有界唯一性已证.
\end{proof}
在半直线 $x>0$ 上, 零 Dirichlet 数据用奇反射, 零 Neumann 用偶反射:
\[
K_{\mathrm D}(t,x,y)=K_D(t,x-y)-K_D(t,x+y),\quad
K_{\mathrm N}(t,x,y)=K_D(t,x-y)+K_D(t,x+y).
\]
积分范围为 $y>0$. 对 $x=0$ 值或导数直接计算得所需边界; 初值由全空间近似单位在内部成立. Dirichlet 核对 $x,y>0$ 非负, 因 $|x-y|<x+y$; Neumann 核总质量为1.
\begin{exercise}
$g=0$, $f(t,x)=1$ 于全空间, 求解并核验初始迹.
\end{exercise}
\begin{solution}
$K_D*1=1$, Duhamel 积分为 $\int_0^t1\dd s=t$. 因而 $u=t$, 有 $u_t=1$, $\Delta u=0$, $u(0)=0$.
\end{solution}

\originalchapter{调和函数与最大值原理}\label{ch:harmonic}
\sourcelecture{6,9}
\originalsection{静态模型与平均值性质}
热方程的无源稳态满足 $\Delta u=0$. 静电场在无旋且单连通的区域内可写 $E=-\nabla u$, 由 $\diver E=\rho$ 得 $-\Delta u=\rho$. 无旋场在任意非单连通区域未必有全局势, 所以不能省略拓扑条件. 本册统一采用 $-\Delta u=f$; 与原课写 $\Delta u=f$ 时的源项相差一个负号.
\begin{definition}
$u\in C^2(\Omega)$ 且 $\Delta u=0$ 时称为调和函数. $\Delta u\ge0$ 称为次调和, $\Delta u\le0$ 称为超调和. 记 $\omega_n=|\Sph^{n-1}|$, 则 $|B_R|=\omega_n R^n/n$.
\end{definition}
\begin{theorem}\label{thm:meanvalue}
若 $\overline B_R(x)\subset\Omega$, $u$ 调和, 则
\[
u(x)=\frac1{\omega_nR^{n-1}}\int_{\partial B_R(x)}u\dd S
=\frac n{\omega_nR^n}\int_{B_R(x)}u\dd y.
\]
次调和函数的中心值不大于这两个平均值.
\end{theorem}
\begin{proof}
设 $A(r)=\omega_n^{-1}\int_{\Sph^{n-1}}u(x+r\theta)\dd\theta$. 链式法则及散度定理给
\[
A'(r)=\frac1{\omega_nr^{n-1}}\int_{\partial B_r(x)}\partial_\nu u\dd S
=\frac1{\omega_nr^{n-1}}\int_{B_r(x)}\Delta u\dd y.
\]
调和时导数为零, $r\downarrow0$ 得 $A(r)=u(x)$. 对球面平均乘 $\omega_nr^{n-1}$ 并从零积分到 $R$, 得球体平均. 次调和时 $A'\ge0$, 从而平均不小于 $A(0)=u(x)$.
\end{proof}
\originalsection{强最大值与边界稳定性}
\begin{theorem}\label{thm:strongharmonic}
连通区域上的调和函数若在内点取得全域最大值或最小值, 则恒为常数. 若 $\Omega$ 有界、$u\in C(\overline\Omega)$, 则最大最小值可在边界取得.
\end{theorem}
\begin{proof}
设最大值为 $M$, $u(x_0)=M$. 任意小球的平均也为 $M$, 而 $M-u\ge0$ 连续, 积分为零, 故小球内 $u=M$. 集合 $\{u=M\}$ 在 $\Omega$ 中既闭又开, 连通性使它等于全域. 最小值对 $-u$ 应用. 有界闭包紧, 极值存在, 若仅在内部取得则常数, 此时边界也取相同值.
\end{proof}
对 $\Delta u\ge0$, 用 $u+\eps|x|^2$ 的严格次调和性排除内部最大, 然后令 $\eps\downarrow0$, 得弱最大值原理. 这里不需假设次调和函数满足平均值等式.
\begin{corollary}\label{cor:poissoncomparison}
若 $-\Delta u=f$, $-\Delta v=g$, 且 $f\le g$, $u\le v$ 于边界, 则 $u\le v$ 于内部. 调和边界问题有至多一个闭包连续解, 并满足 $\|u-v\|_{C(\overline\Omega)}\le\|u-v\|_{C(\partial\Omega)}$.
\end{corollary}
\begin{proof}
$\Delta(u-v)=g-f\ge0$, 对差使用弱最大值原理. 同源项时差调和, 对正负差分别应用得范数估计及唯一性.
\end{proof}
\begin{example}
$\Omega\subset B_R(0)$, $-\Delta u=f$, $u=0$ 于边界. 给出只依赖源上界的估计.
\end{example}
\begin{solution}
令 $M=\|f\|_\infty$, $\psi=M(R^2-|x|^2)/(2n)$. 有 $-\Delta\psi=M$, 且 $\psi\ge0$ 于边界. 比较 $u$ 和 $\psi$, 再比较 $-u$ 和 $\psi$, 得 $|u|\le\psi\le MR^2/(2n)$. 该估计直接体现解对源的连续依赖.
\end{solution}
\originalsection{光滑性与局部估计}
\begin{proposition}\label{prop:harmonicsmooth}
$C^2$ 调和函数实际上为 $C^\infty$. 若 $\overline B_R(x_0)\subset\Omega$, 则
$|\partial^\alpha u(x_0)|\le C_{n,\alpha}R^{-|\alpha|}\sup_{B_R(x_0)}|u|$.
\end{proposition}
\begin{proof}
取径向 $\rho\in C_c^\infty(B_1)$, $\rho\ge0$, $\int\rho=1$, $\rho_r(y)=r^{-n}\rho(y/r)$. 用球面平均和径向积分得 $u*\rho_r=u$ 于距边界大于 $r$ 的地方. 右边卷积光滑, 且 $\partial^\alpha u=u*\partial^\alpha\rho_r$. 取 $r=R/2$, 由 $\|\partial^\alpha\rho_r\|_1=r^{-|\alpha|}\|\partial^\alpha\rho\|_1$ 得估计. 所选光滑核可取 $C\exp[-1/(1-|y|^2)]$ 于单位球内, 外部为零, 常数 $C$ 使积分为1.
\end{proof}
\begin{exercise}
证明全空间有界调和函数为常数.
\end{exercise}
\begin{solution}
固定 $x_0$, 上述一阶导数估计对任意 $R$ 成立, $|\nabla u(x_0)|\le C R^{-1}\|u\|_\infty$. 令 $R\to\infty$ 得 $\nabla u=0$, 连通性给常数. 后章用 Harnack 不等式将“两侧有界”加强为“单侧有界”.
\end{solution}

\originalchapter{Poisson 方程的基本解与弱问题}\label{ch:fundamental}
\sourcelecture{7; 弱存在性为编者补充}
\originalsection{点源与基本解的归一化}
径向函数 $v(r)$ 满足 $\Delta v=v''+(n-1)v'/r$. 令其为零得到 $v'=Cr^{1-n}$, 所以 $n\ge3$ 时 $v=Ar^{2-n}+B$, $n=2$ 时 $v=A\log r+B$. 为使外向通量 $-\int_{\partial B_r}\partial_r\Phi=1$, 取
\begin{equation}\label{eq:newtonkernel}
\Phi(x)=\begin{cases}
\displaystyle\frac{|x|^{2-n}}{(n-2)\omega_n},&n\ge3,\\[3pt]
\displaystyle-\frac1{2\pi}\log|x|,&n=2.
\end{cases}
\end{equation}
本式为 $-\Delta$ 的基本解. 三维为 $1/(4\pi|x|)$. $n-2$ 因子不可省略; 原课一般维数的部分公式缺该因子, 本册按通量核定.
\begin{theorem}\label{thm:fundsolution}
$\Phi$ 局部可积, 且 $-\Delta\Phi=\delta_0$ 于分布意义.
\end{theorem}
\begin{proof}
原点附近 $n\ge3$ 的径向积分为常数乘 $\int_0^\eps r\dd r$, 二维为 $\int_0^\eps r|\log r|\dd r$, 均有限. 给 $\psi\in C_c^\infty$, 取大球包含支撑, 挖去 $B_\eps$. 外边界项为零. 内边界的区域外法向为 $-\partial_r$, 用 Green 第二等式得
\[
\int_{|x|>\eps}\Phi(-\Delta\psi)\dd x
=\int_{|x|=\eps}(\Phi\partial_r\psi-\psi\partial_r\Phi)\dd S.
\]
第一项为 $O(\eps)$, 二维为 $O(\eps|\log\eps|)$, 均趋零. 第二项因 $-\partial_r\Phi=1/(\omega_n\eps^{n-1})$ 等于球面平均 $\psi$, 趋于 $\psi(0)$. 被挖去小球的左端积分也趋零, 故得分布等式.
\end{proof}
\originalsection{Newton 势与无穷远条件}
\begin{theorem}\label{thm:newtonpotential}
$f\in C_c^\infty(\R^n)$ 时 $u=\Phi*f$ 为光滑解 $-\Delta u=f$. $n\ge3$ 时 $u(x)=O(|x|^{2-n})$, 在趋零解中唯一. 二维有
\[
u(x)=-\frac{\int f}{2\pi}\log|x|+O(|x|^{-1}),\qquad |x|\to\infty,
\]
故非零总源时不能同时要求 $u\to0$. 在给定上述渐近式的解中唯一.
\end{theorem}
\begin{proof}
把导数移到紧支撑的 $f$: $\partial^\alpha u(x)=\int\Phi(y)\partial^\alpha f(x-y)\dd y$. 局部可积性使此式有限且连续, 所以 $u$ 光滑. 基本解分布等式给 $-\Delta u=f$, 光滑性使它逐点成立. 若 $\supp f\subset B_a$, $|x|>2a$, 则 $|x-y|\ge|x|/2$, 得高维衰减. 二维写 $\log|x-y|=\log|x|+O(a/|x|)$, 在 $|y|\le a$ 一致成立, 故得展开. 两解之差调和且趋零, 在大球边界差的上界趋零, 最大值原理给全空间差为零.
\end{proof}
若只要求二维 $u/\log|x|\to0$ 与 $\nabla u\to0$, 常数仍满足, 所以该条件不能确定加法常数. 若总源为零, 本册选趋零归一化. 若总源非零, 选上述常数项为零的渐近式.
\begin{example}
三维单位球内均匀源 $f=\rho$, 球外为零, 求趋零的径向弱解.
\end{example}
\begin{solution}
径向方程 $(r^2u')'=-r^2\rho$ 给内部 $u'= -\rho r/3$, $u=A-\rho r^2/6$. 外部为 $B/r$. 因没有表面点源, $u,u'$ 于 $r=1$ 连续, 得 $B=\rho/3$, $A=\rho/2$. 故内部 $u=\rho(3-r^2)/6$, 外部 $u=\rho/(3r)$. 法向导数连续使分部积分无额外界面分布, 确实满足弱方程. 二阶导数于界面跳跃, 与源的跳跃一致.
\end{solution}
\originalsection{齐次 Dirichlet 弱解的存在}
一般区域未必有可写出的 Green 函数. 以下提供存在性的另一个入口, 不把一般光滑域经典解的正则性当作未经证明的结论.
\begin{definition}
$H^1(\Omega)$ 是 $u$ 及全部弱一阶导数在 $L^2$ 中的空间. $H_0^1(\Omega)$ 是 $C_c^\infty(\Omega)$ 在 $H^1$ 范数下的闭包; 齐次边界值在这里由闭包定义表达, 本册不需要一般迹定理.
\end{definition}
\begin{lemma}\label{lem:poincare}
若 $\Omega$ 有界且包含于边长 $L$ 的立方体, 则对 $v\in H_0^1(\Omega)$ 有 $\|v\|_2\le L\|\nabla v\|_2$.
\end{lemma}
\begin{proof}
先令 $v\in C_c^\infty(\Omega)$, 零延拓至立方体. 沿第一坐标 $v(x_1,x')=\int_0^{x_1}\partial_1v(s,x')\dd s$, Cauchy--Schwarz 给 $|v|^2\le L\int_0^L|\partial_1v|^2\dd s$. 再积分 $x_1,x'$ 得 $\|v\|_2^2\le L^2\|\partial_1v\|_2^2$. 对逼近列取极限即可. 因此 $\|\nabla v\|_2$ 是 $H_0^1$ 上与 $H^1$ 范数等价的完备范数.
\end{proof}
\begin{theorem}\label{thm:weakpoisson}
$\Omega$ 有界, $f\in L^2(\Omega)$ 时, 存在唯一 $u\in H_0^1(\Omega)$ 满足
\[
\int_\Omega\nabla u\cdot\nabla v\dd x=\int_\Omega f v\dd x\quad(v\in H_0^1(\Omega)).
\]
并有 $\|\nabla u\|_2\le L\|f\|_2$. 它是泛函 $J(v)=\tfrac12\|\nabla v\|_2^2-\int fv$ 的唯一极小点.
\end{theorem}
\begin{proof}
$H_0^1$ 以梯度内积成为 Hilbert 空间, $\ell(v)=\int fv$ 有界, 因 $|\ell(v)|\le L\|f\|_2\|\nabla v\|_2$. 采用先修的 Hilbert 空间 Riesz 表示定理, 得唯一 $u$ 表示 $\ell$, 正是弱式. 取 $v=u$ 得范数估计. 对 $v=u+w$, 展开并以弱式消去线性项, 得 $J(u+w)=J(u)+\tfrac12\|\nabla w\|_2^2$, 故唯一极小. 唯一性也可对两解之差取自身测试, 不依赖 Riesz 定理.
\end{proof}
这里证明的是能量空间中的弱存在性, 不自动给闭包 $C^2$ 解. 对球域连续边界的调和解, 下一章直接构造并证明取得边界值. 不把所有 $f$ 和所有 Lipschitz 域一概声称具有经典解.

\originalchapter{Green 函数、Poisson 核与 Harnack 不等式}\label{ch:green}
\sourcelecture{7--9}
\originalsection{Green 表示公式}
\begin{definition}
齐次 Dirichlet Green 函数 $G(x,y)$ 满足 $-\Delta_yG(x,y)=\delta_x(y)$, 且边界上 $G(x,\xi)=0$. 一般写 $G(x,y)=\Phi(x-y)-H_x(y)$, 其中 $H_x$ 调和且边界值等于 $\Phi(x-\xi)$. 定义中仍需规定奇异部分及正则性, 不能把任意分布关系当作可逐点使用的核.
\end{definition}
\begin{theorem}\label{thm:greenrepresentation}
若 $\Omega$ 有界光滑, $u\in C^2(\overline\Omega)$, $-\Delta u=f$, $u=g$ 于边界, 且具有上述 Green 函数, 则
\[
u(x)=\int_\Omega G(x,y)f(y)\dd y+\int_{\partial\Omega}P(x,\xi)g(\xi)\dd S_\xi,
\qquad P=-\partial_{\nu_y}G.
\]
\end{theorem}
\begin{proof}
在 $\Omega\setminus\overline B_\eps(x)$ 使用 Green 第二等式, 内部 $\Delta_yG=0$. 小球边界的区域外法向指向 $x$. $G\partial_\nu u$ 项因 $u$ 光滑为 $O(\eps)$, 二维为 $O(\eps|\log\eps|)$; $u\partial_\nu G$ 项趋于 $u(x)$, 因基本解的单位通量. $H_x$ 的平滑部分不贡献极限. 外边界因 $G=0$ 只留下 $-u\partial_\nu G=P g$. 内部 $-\Delta u=f$ 项给 $\int Gf$. 整理即得表示式. 此论证也说明正负号必须与 $-\Delta$ 的约定一致.
\end{proof}
该定理是给定 Green 函数及经典解时的表示, 一般域 Green 函数存在和边界正则性不在这里凭公式推出. 三维半空间 $x_n>0$ 可用像点 $x^*=(x',-x_n)$ 构造 $G=\Phi(x-y)-\Phi(x^*-y)$, 直接满足零边界.
\originalsection{球域像点与 Poisson 公式}
对 $B_R(0)$, $n\ge3$, $x\ne0$ 的反演像点为 $x^*=R^2x/|x|^2$. 若 $|y|=R$, 直接平方可得 $|x^*-y|=(R/|x|)|x-y|$. 因而
\[
G(x,y)=\Phi(x-y)-\left(\frac R{|x|}\right)^{n-2}\Phi(x^*-y).
\]
像点在球外, 第二项于球内调和, 两项在边界抵消. $x=0$ 时取极限 $G(0,y)=\Phi(y)-[(n-2)\omega_nR^{n-2}]^{-1}$. 二维为
\[
G(x,y)=\frac1{2\pi}\log\frac{|x|\,|x^*-y|}{R|x-y|},\qquad
G(0,y)=\frac1{2\pi}\log\frac R{|y|}.
\]
对 $y$ 取外法向导数, 两维与高维统一给
\begin{equation}\label{eq:poissonkernel}
P_R(x,\xi)=\frac{R^2-|x|^2}{\omega_nR|x-\xi|^n},\qquad |\xi|=R.
\end{equation}
例如高维 $\nabla_y\Phi(x-y)=(x-y)/(\omega_n|x-y|^n)$; 第二项同理, 再用像点距离等式及 $\nu_y=\xi/R$ 相减, 分子化为 $|x|^2-R^2$, 最后加负号得\eqref{eq:poissonkernel}.
\begin{theorem}\label{thm:poissonball}
若 $g\in C(\partial B_R)$, 则 $u(x)=\int_{\partial B_R}P_R(x,\xi)g(\xi)\dd S$ 是唯一 $C^\infty(B_R)\cap C(\overline B_R)$ 调和解, 连续取得边界 $g$.
\end{theorem}
\begin{proof}
对固定边界点直接求导得 $\Delta_xP_R(x,\xi)=0$; 可将 $R^2-|x|^2$ 与 $|x-\xi|^{-n}$ 的 Laplace 乘积展开核验, 用 $|\xi|^2=R^2$ 消去全部项. 内部紧集到边界有正距离, 故积分可任意求导. $P_R>0$. 用 Green 表示定理对恒等于1的调和函数, 得 $\int P_R=1$, 该使用没有循环: 恒定解事先已知, 仅借表示式计算核的总量.

对 $x\to\xi_0\in\partial B_R$, 将 $u(x)-g(\xi_0)=\int P_R[g(\xi)-g(\xi_0)]$ 分为近远两段. 近处以 $g$ 的一致连续性控制在 $\eps$; 远处当 $|\xi-\xi_0|\ge\delta$, $x$ 充分近时 $|x-\xi|\ge\delta/2$, 所以核积分至多 $C_\delta(R^2-|x|^2)\to0$. 得连续边界迹. 唯一性由最大值原理.
\end{proof}
\begin{example}
单位圆上 $g(\theta)=\cos(3\theta)+2\sin(2\theta)$, 求调和延拓.
\end{example}
\begin{solution}
极坐标 Laplace 算子为 $\partial_{rr}+r^{-1}\partial_r+r^{-2}\partial_{\theta\theta}$. 逐项代入得 $r^3\cos3\theta$ 和 $r^2\sin2\theta$ 调和. 它们在原点也是多项式: 分别为 $x^3-3xy^2$ 和 $2xy$. 因而 $u=r^3\cos3\theta+2r^2\sin2\theta$ 连续取得边界值, 唯一性认定它就是 Poisson 积分.
\end{solution}
\originalsection{Harnack 与单侧 Liouville 定理}
\begin{theorem}\label{thm:harnack}
若 $u\ge0$ 在 $B_R(0)$ 内调和, 则对 $r=|x|<R$,
\[
\frac{R^{n-2}(R-r)}{(R+r)^{n-1}}u(0)
\le u(x)\le
\frac{R^{n-2}(R+r)}{(R-r)^{n-1}}u(0).
\]
\end{theorem}
\begin{proof}
先对球半径 $s<R$ 且 $s>r$ 使用 Poisson 公式, 边界值非负. $s-r\le|x-\xi|\le s+r$ 给上下核界, 而球面平均给 $\int_{\partial B_s}u=\omega_n s^{n-1}u(0)$. 相乘得到相同公式中 $R$ 换成 $s$. 令 $s\uparrow R$ 即得, 不要求 $u$ 连续延伸到最外球面.
\end{proof}
\begin{corollary}
若全空间调和函数有统一上界或下界, 则为常数. 非负调和函数若在内点为零, 则在连通域恒零.
\end{corollary}
\begin{proof}
下界 $M$ 时令 $v=u-M\ge0$, 对任意固定 $x$ 在任意大球用 Harnack, 两侧系数均趋1, 得 $v(x)=v(0)$. 上界对 $-u$ 使用. 零点情形在小球由 Harnack 得恒零, 再以连通性沿相交球传播, 或直接用强最小值原理.
\end{proof}
\begin{exercise}
若三维 $-\Delta u=f\ge0$, $u=0$ 于光滑有界域边界, 证明 $u\ge0$, 并说明正 Green 函数与此结论的一致性.
\end{exercise}
\begin{solution}
$\Delta(-u)=f\ge0$, 弱最大值原理给 $-u\le0$. Green 函数固定点源的正性可在挖去小球后比较: 小球边界 $\Phi-H_x>0$ 对充分小半径成立, 外边界零, 中间调和, 故 $G\ge0$. 因此表示 $u=\int Gf$ 也给非负性.
\end{solution}

\originalchapter{波方程与显式传播公式}\label{ch:wave}
\sourcelecture{10--12}
\originalsection{声学线性化与一维特征}
一维可压缩流体的连续方程和动量方程为 $\rho_t+(\rho v)_x=0$, $(\rho v)_t+(\rho v^2+p)_x=0$. 在静态 $\rho=\rho_0>0$, $v=0$ 附近取 $\rho=\rho_0+\eta$ 并保留一阶小量, 得
\[
\eta_t+\rho_0v_x=0,\qquad \rho_0v_t+p'(\rho_0)\eta_x=0.
\]
对第一式求时间导数, 用第二式的空间导数消去 $v$, 得 $\eta_{tt}-c^2\eta_{xx}=0$, $c^2=p'(\rho_0)>0$. 若 $p=K\rho^\gamma$, 则 $c^2=K\gamma\rho_0^{\gamma-1}$. 线性波方程是小扰动近似, 不是完整非线性流体方程.
\begin{theorem}[d'Alembert 公式]\label{thm:dalembert}
$f\in C^2(\R)$, $g\in C^1(\R)$, $c>0$ 时, $u_{tt}-c^2u_{xx}=0$, $u(0)=f$, $u_t(0)=g$ 的唯一 $C^2$ 解为
\begin{equation}\label{eq:dalembert}
u(t,x)=\frac{f(x-ct)+f(x+ct)}2+\frac1{2c}\int_{x-ct}^{x+ct}g(y)\dd y.
\end{equation}
\end{theorem}
\begin{proof}
取 $\xi=x+ct$, $\eta=x-ct$, 则 $\partial_t^2-c^2\partial_x^2=-4c^2\partial_\xi\partial_\eta$. 因而全部 $C^2$ 解为 $A(\xi)+B(\eta)$. 初时 $A+B=f$, $c(A'-B')=g$, 得 $A'=(f'+g/c)/2$, $B'=(f'-g/c)/2$. 从 $x$ 分别积分至 $x+ct,x-ct$, 并用 $A(x)+B(x)=f(x)$, 即得公式. 反过来公式的导数给方程, $t=0$ 得初值, $u_t(0)=g$. 以上从任意解推导也给唯一性, 不需全空间衰减假设.
\end{proof}
右行波为 $F(x-ct)$, 左行波为 $G(x+ct)$. 速度初值产生区间积分, 一般解不只是一个刚性平移. 点 $(t,x)$ 只依赖区间 $[x-ct,x+ct]$ 内的数据.
\begin{center}
\begin{tikzpicture}[x=1.3cm,y=1.05cm,>=Latex]
\draw[->](-2.5,0)--(2.6,0) node[right]{$x$};
\draw[->](0,-.1)--(0,2.5) node[above]{$t$};
\draw[structurecolor,thick](-2,0)--(0,2)--(2,0);
\draw[main,thick](-2,0)--(2,0);
\fill(0,2)circle(1.5pt) node[right=4pt]{$(t_0,x_0)$};
\node[below] at (-2,0){$x_0-ct_0$};\node[below] at (2,0){$x_0+ct_0$};
\node at(0,0.7){依赖区域};
\end{tikzpicture}
\end{center}
图为一维后向特征锥的示意, 横坐标原点平移到 $x_0$, 斜率与坐标尺度对应.
\originalsection{反射与有界弦}
半直线固定端 $u(t,0)=0$ 用 $f,g$ 的奇延拓套用\eqref{eq:dalembert}. 若要求闭半平面的 $C^2$ 解, 奇延拓 $f$ 需 $f(0)=f''(0)=0$, 奇延拓 $g$ 需 $g(0)=0$; 仅 $f(0)=g(0)=0$ 不保证所有二阶导数跨反射线连续. 弱解可用较弱条件. 自由端 $u_x(t,0)=0$ 用偶延拓, 对 $C^2$ 构造需 $f'(0)=g'(0)=0$.

固定端区间 $0<x<L$ 的模态为
\[
u(t,x)=\sum_{k\ge1}\left(a_k\cos(c\sqrt{\lambda_k}t)+\frac{b_k}{c\sqrt{\lambda_k}}\sin(c\sqrt{\lambda_k}t)\right)\sin(k\pi x/L),
\]
$a_k,b_k$ 分别为位移和速度的正弦系数. 对足够光滑相容数据可逐项求导; 对 $f\in H_0^1$, $g\in L^2$, 级数在能量范数中收敛. 这里的能量收敛可由导数的 Parseval 补足: 对光滑紧支撑 $f$, 分部积分给 $f'$ 的余弦系数为 $(k\pi/L)a_k$, 常数系数为零. 偶延拓与 Fejér 核的论证给余弦系完备性, 故 $\|f'\|_2^2=(L/2)\sum_k\lambda_k|a_k|^2$. 用 $C_c^\infty$ 在 $H_0^1$ 中的定义性稠密推广到一般 $f$. 速度有 $(L/2)\sum_k|b_k|^2=\|g\|_2^2$. 每个模态的 $c^2\lambda_k$ 倍位移平方与速度平方之和恒为 $c^2\lambda_k|a_k|^2+|b_k|^2$, 可和尾项一致控制所有时间的能量误差. 因此级数在 $H_0^1\times L^2$ 中收敛, 时间连续由有限模态近似和统一尾界得到. 各模态能量恒定, 不像热模态指数衰减.
\begin{example}
$f(x)=0$, $g(x)=1$ 于全直线, 求波方程解并解释其不是单个行波.
\end{example}
\begin{solution}
\eqref{eq:dalembert}给 $u=(2c)^{-1}(2ct)=t$. 它在每个空间点同时线性增长, $u_{tt}=u_{xx}=0$. 可分成两个线性行波的叠加, 但其整体图形不是同一轮廓的单向平移.
\end{solution}
\originalsection{球面平均与 Kirchhoff 公式}
以下先令 $c=1$. 对三维函数 $h$ 定义 $M_rh(x)=(4\pi)^{-1}\int_{\Sph^2}h(x+r\theta)\dd\theta$. 散度定理给
\[
\partial_r M_rh=\frac1{4\pi r^2}\int_{B_r(x)}\Delta h\dd y,
\qquad \partial_r(r^2\partial_rM_rh)=r^2M_r\Delta h.
\]
因此 $U(t,r)=M_ru(t,\cdot)(x)$ 满足 $U_{tt}=U_{rr}+2U_r/r$, 而 $V=rU$ 满足 $V_{tt}=V_{rr}$ 且 $V(t,0)=0$. 其初值为 $rM_rf$, $rM_rg$, 皆对 $r$ 作奇延拓. 用一维公式后除以 $r$ 并令 $r\downarrow0$, 得
\begin{equation}\label{eq:kirchhoff}
u(t,x)=\partial_t\bigl(tM_tf(x)\bigr)+tM_tg(x).
\end{equation}
\begin{theorem}\label{thm:kirchhoff}
$f\in C^3(\R^3)$, $g\in C^2(\R^3)$ 时, \eqref{eq:kirchhoff}给三维 $C^2$ 初值解. 对速度 $c$, 公式为 $u=\partial_t(tM_{ct}f)+tM_{ct}g$.
\end{theorem}
\begin{proof}
前面的推导表明任意经典解必须取此形式. 还需核验构造确实满足方程. $M_th$ 满足 Darboux 恒等式 $(\partial_t^2+2t^{-1}\partial_t)M_th=M_t\Delta h$. 故 $W_h=tM_th$ 满足 $(W_h)_{tt}=\Delta_xW_h$. 时间导数 $\partial_tW_f$ 也满足波方程. $W_h(0)=0$, $W_h'(0)=h$, $W_h''(0)=0$, 后一式由球面对 $\theta$ 的平均为零及 Taylor 展开得到. 所以 $u(0)=f$, $u_t(0)=g$.

$f$ 三阶、$g$ 二阶连续保证公式中所需的两阶导数存在且连续, 正时间仅在紧球面求导; 在 $t=0$ 用 $M_th=h+t^2\Delta h/6+o(t^2)$ 及相应导数展开延伸. 一般速度由时间缩放或直接对 $M_{ct}$ 的恒等式得到. 唯一性亦由下一章的局部能量法保证.
\end{proof}
三维公式只在球面 $|y-x|=ct$ 取初值及其邻域导数, 这是尖锐 Huygens 原理. 若初值紧支撑, 某固定点最终不再受其影响. 一维速度积分则保留锥内部贡献, 没有同样的尖锐性质.
\originalsection{二维降维与源项}
将二维数据视为三维中与 $x_3$ 无关的函数. 上下半球投影到二维圆盘, 球面元为 $t\,\dd y/\sqrt{t^2-|x-y|^2}$, 两半球相加, 因而
\[
tM_th(x)=\frac1{2\pi}\int_{|y-x|<t}\frac{h(y)}{\sqrt{t^2-|x-y|^2}}\dd y.
\]
代入\eqref{eq:kirchhoff}得到二维 Poisson 波公式. 积分核在圆周奇异但可积, 初值依赖整个圆盘, 所以没有三维的尖锐 Huygens 性质. 正则性和 PDE 可由三维延拓公式核验, 避免直接在移动奇异边界上形式求导.

记 $S_c(t)h$ 为零位移、速度 $h$ 的齐次波解. 对光滑源 $F$, Duhamel 公式为
\[
u(t)=u_{\mathrm{hom}}(t)+\int_0^tS_c(t-s)F(s)\dd s.
\]
因为 $S_c(0)=0$, $\partial_tS_c(0)=I$, 两次时间求导给源项 $F(t)$ 加上齐次波算子作用, 故 $(\partial_t^2-c^2\Delta)u=F$. 这与热方程的一次时间求导有不同的端点项.
\begin{exercise}
三维 $f=0$, $g=1$, 证明 $u=t$ 并核对一般速度公式中的系数.
\end{exercise}
\begin{solution}
球面平均 $M_{ct}1=1$, 所以 $u=t$. 若改写成球面积分, 速度项应是 $(4\pi c^2t)^{-1}\int_{|y-x|=ct}g\dd S$, 球面积为 $4\pi c^2t^2$, 恰得 $t$. 漏掉 $c^2$ 会使初始速度错误.
\end{solution}

\originalchapter{波能量与 Minkowski 几何}\label{ch:energywave}
\sourcelecture{12--14}
\originalsection{守恒与锥上的能量}
设 $u_{tt}-c^2\Delta u=F$. 定义密度 $e=\tfrac12(u_t^2+c^2|\nabla u|^2)$, 通量 $p=-c^2u_t\nabla u$. 直接求导得
\begin{equation}\label{eq:localwaveenergy}
\partial_te+\diver p=u_tF.
\end{equation}
这条局部等式包含整体守恒及有限传播. 齐次 Dirichlet 边界的光滑解有 $u_t=0$, 齐次 Neumann 有 $\partial_\nu u=0$, 都使边界通量消失; Robin $\partial_\nu u+\alpha u=0$, $\alpha=\alpha(x)\ge0$ 与时间无关, 则能量要加 $\frac{c^2}2\int_{\partial\Omega}\alpha u^2$ 才守恒.
\begin{theorem}\label{thm:coneenergy}
设 $u$ 为齐次波方程的 $C^2$ 解, $0\le t<R/c$. 对收缩球 $B_{R-ct}(x_0)$ 定义 $E(t)=\int e(t,x)\dd x$, 则 $E(t)\le E(0)$. 相同初值数据在 $B_R(x_0)$ 上的两解在后向锥 $|x-x_0|<R-ct$ 内相同.
\end{theorem}
\begin{proof}
对移动球积分求导, 边界向内速度为 $c$, 所以
\begin{align*}
E'(t)&=\int_{\partial B_{R-ct}}\bigl(c^2u_t\partial_\nu u-ce\bigr)\dd S\\
&=-\frac c2\int_{\partial B_{R-ct}}
\left[(u_t-c\partial_\nu u)^2+c^2|\nabla_{\mathrm{tan}}u|^2\right]\dd S\le0.
\end{align*}
$\nabla_{\mathrm{tan}}$ 为球面切向梯度, $|\nabla u|^2=(\partial_\nu u)^2+|\nabla_{\mathrm{tan}}u|^2$. 对两解差 $w$, 初始能量零, 所以锥内 $w_t=\nabla w=0$. 沿固定空间点的竖直线连接初始截面, 因该线仍在锥内, 且 $w(0,x)=0$, 得 $w(t,x)=0$. 仅控制导数后还需初始位移去掉常数.
\end{proof}
本证明在紧锥上进行, 不要求全空间总能量有限. 因此任意 $C^2$ 全空间波初值解均唯一. 初值的支撑包含于 $B_R$ 时, 支撑在 $t$ 时刻包含于 $B_{R+ct}$, 由外部各后向锥零数据得到.
\begin{proposition}\label{prop:wavestability}
在消失通量的边界条件下或全空间充分衰减时,
\[
E(t)^{1/2}\le E(0)^{1/2}+\frac1{\sqrt2}\int_0^t\|F(s)\|_2\dd s.
\]
\end{proposition}
\begin{proof}
积分\eqref{eq:localwaveenergy}得 $E'=\int u_tF\le\sqrt{2E}\|F\|_2$. 用 $(E+\eps^2)^{1/2}$ 处理零点, 积分并令 $\eps\downarrow0$. 这给能量空间的连续依赖. $E$ 不直接控制 $u$ 的 $L^2$ 范数, 但 $\|u(t)\|_2\le\|f\|_2+\int_0^t\|u_t(s)\|_2\dd s$ 可补充该控制.
\end{proof}
\originalsection{Lorentz 变换与零标架}
以下速度归一化为1. 时空坐标 $x^0=t,x^1,\ldots,x^n$, 度量矩阵 $m=\operatorname{diag}(-1,1,\ldots,1)$. $m(X,Y)=-X^0Y^0+\sum_iX^iY^i$. 非零向量满足 $m(X,X)<0$、$=0$、$>0$ 时分别称为类时、零、类空. 因 $m(X,X)$ 可负, 它不是正定内积. $X^0>0$ 称为未来指向.

指标升降用 $X_\mu=m_{\mu\nu}X^\nu$, $X^\mu=m^{\mu\nu}X_\nu$; 重复的一上一下指标对 $0,\ldots,n$ 求和. $\partial^0=-\partial_t$, $\partial^i=\partial_i$, $\square=m^{\mu\nu}\partial_\mu\partial_\nu=-\partial_t^2+\Delta$.
\begin{proposition}
若线性变换 $\Lambda$ 满足 $\Lambda^Tm\Lambda=m$, 则保持类时/零/类空类型, 且 $u\circ\Lambda$ 与 $u$ 的波算子满足 $\square(u\circ\Lambda)=(\square u)\circ\Lambda$.
\end{proposition}
\begin{proof}
第一项直接由 $m(\Lambda X,\Lambda X)=m(X,X)$. 取逆矩阵得等价关系 $\Lambda m^{-1}\Lambda^T=m^{-1}$, 链式法则使二阶导数系数化为同一个 $m^{-1}$, 故得算子等式. 取行列式得 $|\det\Lambda|=1$. 这只保证度量保持, 不自动保证未来指向保持; 需另选保时间方向的分支.
\end{proof}
一维 Lorentz boost 为 $(t',x')=(\gamma(t-vx),\gamma(x-vt))$, $|v|<1$, $\gamma=(1-v^2)^{-1/2}$. 平方展开验证 $-(t')^2+(x')^2=-t^2+x^2$.

令 $r=|x|>0$, $L=\partial_t+\partial_r$, $\underline L=\partial_t-\partial_r$, 配球面上的局部正交切向标架 $e_1,\ldots,e_{n-1}$. 有 $m(L,\underline L)=-2$, $m(L,L)=m(\underline L,\underline L)=0$, 且
\[
m^{-1}=-\tfrac12(L\otimes\underline L+\underline L\otimes L)+\sum_A e_A\otimes e_A.
\]
对基向量逐对求值即可核验. 球面切向标架只需局部选取, 不声称所有维数的球面都可全局平滑选一组标架.
\originalsection{应力张量与乘子}
\begin{definition}
对实标量场 $u$, 令 $Q=m^{\alpha\beta}\partial_\alpha u\partial_\beta u=-u_t^2+|\nabla u|^2$. 能量动量张量为
\[
T^{\mu\nu}=\partial^\mu u\partial^\nu u-\tfrac12m^{\mu\nu}Q.
\]
对向量场 $X$ 的流为 $J_X^\mu=T^{\mu\nu}X_\nu$, 形变张量为 $\pi^X_{\mu\nu}=\tfrac12(\partial_\mu X_\nu+\partial_\nu X_\mu)$.
\end{definition}
\begin{proposition}\label{prop:stress}
$\partial_\mu T^{\mu\nu}=(\square u)\partial^\nu u$. 若 $\square u=0$, 则 $\partial_\mu J_X^\mu=T^{\mu\nu}\pi^X_{\mu\nu}$. 特别地, Killing 场 $\pi^X=0$ 产生守恒流.
\end{proposition}
\begin{proof}
乘积法则得 $(\square u)\partial^\nu u+\partial^\mu u\partial_\mu\partial^\nu u-\tfrac12\partial^\nu Q$. 最后两项因混合导数交换抵消. 再对 $T^{\mu\nu}X_\nu$ 求导, $T$ 对称使反对称部分不贡献, 剩余为 $T:\pi$.
\end{proof}
取 $X=-\partial_t$, $X_0=1$, 得 $J^0=T^{00}=\tfrac12(u_t^2+|\nabla u|^2)$, $J^i=-u_tu_i$, 就是局部能量流. 指标符号若改成 $X=\partial_t$, 流整体变号, 不能忽略.
\begin{lemma}[主导能量正性]\label{lem:dominant}
这里因果向量指类时或零向量, 包括零向量; 非零未来指向表示时间分量为正. 若 $X,Y$ 同为未来指向因果向量, 则 $T_{\mu\nu}X^\mu Y^\nu\ge0$; 同为过去指向也成立.
\end{lemma}
\begin{proof}
先令 $X$ 类时, 可用保未来的 Lorentz 变换把它化为 $(a,0,\ldots,0)$, $a>0$: 沿 $X$ 空间方向用速度 $v=|X'|/X^0<1$ 的 boost 消去空间分量. 张量表达在这种坐标变化下保持. 因 $Y^0\ge|Y'|$, 得
\[
T(X,Y)=a\left[\tfrac12(u_t^2+|\nabla u|^2)Y^0+u_t\nabla u\cdot Y'\right]
\ge\frac{aY^0}2(|u_t|-|\nabla u|)^2\ge0.
\]
类时近似零向量并取极限得零情形. 过去指向同时乘 $-1$ 不改变双线性值.
\end{proof}
注意在 $1+n$ 维中 $m_{\mu\nu}T^{\mu\nu}=(1-(n+1)/2)Q$. 三维空间即 $-Q$, 一般不为零. 后面的课程作业勘误与 Morawetz 修正流都依赖这一点.

\originalchapter{主部分类与 Fourier 变换}\label{ch:fourier}
\sourcelecture{15--18}
\originalsection{二次型与方程类型}
常系数二阶算子写作 $Lu=\sum_{j,k=0}^n a^{jk}u_{x_jx_k}+\sum_jb^ju_{x_j}+du$, 混合导数交换允许将 $A=(a^{jk})$ 对称化. 主符号为 $p(\xi)=\xi^TA\xi$. 非奇异线性坐标变换 $y=Mx$ 使主矩阵成为 $MAM^T$, 是合同变换.
\begin{theorem}\label{thm:classification}
对称常矩阵 $A$ 可经非奇异线性变换化为只有 $1,-1,0$ 的对角阵. 全部特征值同号且非零为椭圆型; 恰一个异号而其余同号且均非零为波型双曲型; 恰一个零而其余同号非零为退化抛物主部. 各种正、负、零方向的数量在非奇异坐标变换下不变.
\end{theorem}
\begin{proof}
由实对称谱定理选正交 $Q$, $QAQ^T=\operatorname{diag}(\lambda_j)$. 非零方向再缩放 $|\lambda_j|^{-1/2}$, 零方向保留尺度, 得所示正规形. 惯性不变可从二次型正定子空间的最大维数证明: 可逆变换双射正定子空间, 故最大维数不变; 负定同理, 零维数由秩决定. 因此类型不依赖所选坐标.
\end{proof}
主部退化尚不足以保证有时间演化. 如 $u_{xx}=0$ 在 $(t,x)$ 上主部秩1, 但没有 $u_t$; 热型演化还要求低阶项在主部零方向有非零一阶导数, 其符号决定时间方向. 变量系数问题的点态分类也不自动给全局适定性.
\begin{example}
把 $u_{tt}-4u_{tx}+2u_{xx}=0$ 化为标准波方程.
\end{example}
\begin{solution}
令 $s=t$, $y=x+2t$. 则 $\partial_t=\partial_s+2\partial_y$, $\partial_x=\partial_y$, 主部为 $\partial_s^2-2\partial_y^2$. 故为速度 $\sqrt2$ 的波方程. 原始混合项不是新的非线性, 只反映坐标选择.
\end{solution}
\originalsection{变换、衰减与微分}
\begin{definition}
对 $f\in L^1(\R^n)$, 定义 $\widehat f(\xi)=\int f(x)\ee^{-2\pi\ii x\cdot\xi}\dd x$, 逆变换为 $h^\vee(x)=\int h(\xi)\ee^{2\pi\ii x\cdot\xi}\dd\xi$. Schwartz 空间 $\Sch$ 由所有 $C^\infty$ 函数组成, 满足每个 $\alpha,\beta$ 的 $\sup_x|x^\alpha\partial^\beta f(x)|<\infty$.
\end{definition}
$L^1$ 变换有界连续, $\|\widehat f\|_\infty\le\|f\|_1$: 绝对值估计给有界性, 用 $2|f|$ 控制指数差, 控制收敛给连续性. 一致连续性可先把 $f$ 尾积分截小, 有界区域内用指数的 Lipschitz 估计.
\begin{proposition}\label{prop:fourierprops}
对 $f,g\in\Sch$, 有
\begin{align*}
\widehat{f(\cdot-a)}(\xi)&=\ee^{-2\pi\ii a\cdot\xi}\widehat f(\xi),&
\widehat{\ee^{2\pi\ii b\cdot x}f}(\xi)&=\widehat f(\xi-b),\\
\widehat{f(x/\lambda)}(\xi)&=|\lambda|^n\widehat f(\lambda\xi)\quad(\lambda\ne0),&
\widehat{f*g}&=\widehat f\widehat g,\\
\widehat{\partial^\alpha f}&=(2\pi\ii\xi)^\alpha\widehat f,&
\partial_\xi^\beta\widehat f&=\widehat{(-2\pi\ii x)^\beta f}.
\end{align*}
此外 $\widehat f\in\Sch$.
\end{proposition}
\begin{proof}
前三式分别换元 $y=x-a$, 合并指数, 换元 $y=x/\lambda$, Jacobian 必须取 $|\lambda|^n$. 卷积式对绝对可积双积分使用 Fubini, 置 $z=x-y$ 后指数分离. 导数式对空间积分分部, Schwartz 衰减消去端点; 频率导数式在积分内求导, 因 $x^\beta f$ 可积. 对频率导数再乘任意频率单项式, 上述两式将它化为某个 Schwartz 函数的 Fourier 变换, 其一致上界由 $L^1$ 范数控制, 所以变换仍在 $\Sch$.
\end{proof}
$L^1$ 函数的变换还在无穷远趋零: 先对紧支撑光滑函数用积分分部得到衰减, 再以其在 $L^1$ 中的稠密性和一致变换估计逼近. 有界连续不等于无穷远取得最大值, 所以范数一般写上确界.
\originalsection{Gaussian 与反演}
\begin{lemma}\label{lem:complexgaussian}
$\operatorname{Re}z>0$ 时,
\[
\widehat{\ee^{-\pi z|x|^2}}(\xi)=z^{-n/2}\ee^{-\pi|\xi|^2/z}.
\]
右半平面平方根取正实轴上正值的连续分支.
\end{lemma}
\begin{proof}
一维令 $I(z)=\int\ee^{-\pi z x^2}\dd x$. 在右半平面紧集内可对 $z$ 求导, 控制函数为多项式乘衰减 Gaussian. 对 $(x\ee^{-\pi z x^2})'$ 积分得 $I=-2zI'$, 所以 $z^{1/2}I$ 的导数为零, 由 $I(1)=1$ 得 $I=z^{-1/2}$. 设 $H(\xi)$ 为变换. 频率求导并以 $(\ee^{-\pi zx^2})'=-2\pi zx\ee^{-\pi zx^2}$ 分部积分, 得 $H'=-2\pi\xi H/z$. 此线性 ODE 与 $H(0)=I(z)$ 给公式, 不需对可能为零的 $H$ 取对数. 高维由 Fubini 和各坐标乘积得到.
\end{proof}
\begin{theorem}[Fourier 反演]\label{thm:inversion}
$f\in\Sch$ 时 $f=(\widehat f)^\vee$. 更一般地, 若 $f$ 连续、$f\in L^1$ 且 $\widehat f\in L^1$, 则逐点反演成立.
\end{theorem}
\begin{proof}
直接交换无正则化的双积分不合法. 对 $\eps>0$ 插入 $\ee^{-\pi\eps|\xi|^2}$, Gaussian 引理及 Fubini 得
\[
\int\widehat f(\xi)\ee^{2\pi\ii x\cdot\xi}\ee^{-\pi\eps|\xi|^2}\dd\xi
=\int f(y)\eps^{-n/2}\ee^{-\pi|x-y|^2/\eps}\dd y.
\]
右边是近似单位, 在连续点趋 $f(x)$. 对连续 $L^1$ 函数, 近点连续性控制, 远点用 $L^1$ 尾与 Gaussian 在远点的上界控制, 不要求 $f$ 全局有界. 左边由 $|\widehat f|$ 控制趋逆变换. 得反演式; 换指数符号给另一个逆次序.
\end{proof}
\originalsection{Plancherel 与 PDE 乘子}
\begin{theorem}\label{thm:plancherel}
Fourier 变换在 $\Sch$ 上保复内积 $\langle f,g\rangle=\int f\overline g$, 并唯一延拓为 $L^2(\R^n)$ 上的酉算子, 逆为逆变换的 $L^2$ 延拓.
\end{theorem}
\begin{proof}
$\Sch$ 上 Fubini 合法, 由反演 $\overline{g(x)}=\int\overline{\widehat g(\xi)}\ee^{-2\pi\ii x\cdot\xi}\dd\xi$, 得
\[
\int f\overline g\dd x=\int\widehat f\,\overline{\widehat g}\dd\xi.
\]
$\Sch$ 在 $L^2$ 中稠密: 先截断空间和函数值, 再用光滑近似单位平滑, 必要时乘紧支撑截断, 误差由平移连续性及尾积分控制. 给 $f\in L^2$ 取 $f_j\in\Sch$ 逼近, 保范数使 $\widehat f_j$ Cauchy, 完备性给极限且与逼近列无关. 逆变换同样延拓, $\Sch$ 上两者互逆, 由连续性在 $L^2$ 上互逆, 故为酉算子.
\end{proof}
$L^2$ 变换一般是范数极限, 不保证原定义积分绝对收敛. 这一点使能量解可以严格处理.

热方程取空间变换得到 $\partial_t\widehat u=-4\pi^2D|\xi|^2\widehat u$, 所以 $\widehat u=\ee^{-4\pi^2Dt|\xi|^2}\widehat g$, Gaussian 反演回到热核. 波方程得到 $\widehat u_{tt}+(2\pi c|\xi|)^2\widehat u=0$, 故
\[
\widehat u(t,\xi)=\cos(2\pi c|\xi|t)\widehat f(\xi)
+\frac{\sin(2\pi c|\xi|t)}{2\pi c|\xi|}\widehat g(\xi).
\]
零频处第二因子定义为 $t$, 是可去奇点. $f\in H^1$, $g\in L^2$ 时乘子有界控制 $u\in C(H^1)$, $u_t\in C(L^2)$, 控制收敛给时间连续及初值. 对光滑近似列传递方程和能量守恒, 得弱能量解. 唯一性可在时间上正则化差解作能量法, 或在变换端解每个频率的 ODE; 这明确了全空间弱解的类别.
\begin{exercise}
计算 $f(x)=\ee^{-\pi a|x|^2}$ 的 $L^2$ 范数与变换范数, $a>0$.
\end{exercise}
\begin{solution}
$\|f\|_2^2=(2a)^{-n/2}$. 变换为 $a^{-n/2}\ee^{-\pi|\xi|^2/a}$, 平方积分为 $a^{-n}(a/2)^{n/2}=(2a)^{-n/2}$, 与 Plancherel 一致.
\end{solution}

\originalchapter{Schrödinger 方程与色散}\label{ch:schrodinger}
\sourcelecture{19--20}
\originalsection{频率演化与解类别}
本章使用原课归一化 $\ii u_t+\tfrac12\Delta u=0$. 复数未知函数同时编码振幅与相位. 平面波 $\ee^{\ii(k\cdot x-\omega t)}$ 代入得 $\omega=|k|^2/2$, 相速度为 $|k|/2$, 群速度为 $\nabla_k\omega=k$. 相速度与波包的群速度不同; 两者均随频率变化, 这是色散.
\begin{theorem}\label{thm:schroL2}
对 $g\in L^2(\R^n)$, 定义
\[
U(t)g=\mathcal F^{-1}\left(\ee^{-2\pi^2\ii t|\xi|^2}\widehat g(\xi)\right),\qquad t\in\R.
\]
则 $U(t)$ 是强连续酉群, 给唯一 $C(\R;L^2)$ 的分布初值解. $g\in\Sch$ 时为经典解并保持 $\Sch$; $H^s$ 正则性保持, 不产生热方程那种一般平滑效应.
\end{theorem}
\begin{proof}
乘子模为1, Plancherel 给保范数、群性质 $U(t+s)=U(t)U(s)$ 及逆 $U(-t)$. 时间连续由 $|\ee^{-2\pi^2\ii t|\xi|^2}-1|^2|\widehat g|^2\le4|\widehat g|^2$ 和控制收敛. 对 $\Sch$ 数据可任意求导, 变换 ODE 为 $\partial_t\widehat u=-2\pi^2\ii|\xi|^2\widehat u$, 等价于 PDE. 对一般数据用 $\Sch$ 逼近, 在测试函数意义传递导数, 得分布解.

说明唯一性时不能随意对一个仅为 $L^2$ 的函数逐点求频率导数. 对频率乘光滑紧支撑截断 $\chi_R$, PDE 给 $v_t=-2\pi^2\ii|\xi|^2v$ 于分布时间意义, 右端为连续 $L^2$ 函数, 所以 $v$ 可微并满足 Hilbert 空间 ODE. 乘 $\ee^{2\pi^2\ii t|\xi|^2}$ 得常数, 初值决定 $v$. 令截断趋1给全部频率唯一. $H^s$ 定义为加权范数 $\int(1+|\xi|^2)^s|\widehat g|^2<\infty$, 模1保持该范数.
\end{proof}
\originalsection{振荡核与色散估计}
\begin{theorem}\label{thm:schrokernel}
$t>0$, $g\in\Sch$ 时
\[
U(t)g(x)=\int_{\R^n}K(t,x-y)g(y)\dd y,
\qquad K(t,x)=(2\pi\ii t)^{-n/2}\ee^{\ii|x|^2/(2t)}.
\]
$\ii^{1/2}=\ee^{\ii\pi/4}$. 对 $g\in L^1\cap L^2$ 有
\[
\|U(t)g\|_\infty\le(2\pi |t|)^{-n/2}\|g\|_1,\qquad t\ne0.
\]
\end{theorem}
\begin{proof}
在频率乘子加 $\ee^{-2\pi^2\delta t|\xi|^2}$, $\delta>0$, 复杂 Gaussian 公式给核
\[
K_\delta(t,x)=[2\pi t(\delta+\ii)]^{-n/2}
\exp\left(-\frac{|x|^2}{2t(\delta+\ii)}\right).
\]
当 $\delta\downarrow0$, 频率端积分由 $|\widehat g|$ 控制, 趋 $U(t)g$. 空间端对固定 $t>0$ 有 $|K_\delta|\le(2\pi t)^{-n/2}$, 故由 $|g|$ 控制, 趋 $K*g$. 这证明核式, 没有将不可积的 $K$ 直接代入仅适用于 $L^1$ 核的卷积变换定理. 取绝对值给色散估计. $t<0$ 由共轭核和群的逆得到. 一般 $L^1\cap L^2$ 数据同时在两范数中逼近: 空间卷积一致收敛, $L^2$ 解收敛, 所以两者几乎处处相同, 估计传递.
\end{proof}
$|K(t,x)|$ 与 $x$ 无关, 因而 $K(t,\cdot)\notin L^1$. 核不是非负平均, 色散衰减由振荡结构产生, 不能使用热方程比较原理. 对 $g\in\Sch$, 初始点态极限无需本册之外的驻相法: 频率积分中 $\widehat g\in L^1$, 直接控制收敛给 $U(t)g(x)\to g(x)$, 甚至一致收敛.
\begin{example}
初值 $g(x)=\ee^{-a|x|^2}$, $a>0$, 求自由演化.
\end{example}
\begin{solution}
将热 Gaussian 解的扩散时间替换为 $\ii t/2$, 或用 Fourier Gaussian 积分, 得
\[
u(t,x)=(1+2\ii at)^{-n/2}\exp\left(-\frac{a|x|^2}{1+2\ii at}\right).
\]
分支从 $t=0$ 连续选取. 模平方为 $(1+4a^2t^2)^{-n/2}\exp[-2a|x|^2/(1+4a^2t^2)]$, 积分为 $(\pi/(2a))^{n/2}$ 不变. 峰值下降与总概率守恒可以同时成立.
\end{solution}
\originalsection{守恒律、势与 Duhamel 公式}
对 $\Sch$ 解, $u_t=\ii\Delta u/2$, 乘 $\overline u$ 积分取实部, 得
\[
\frac{\mathrm d}{\mathrm dt}\int|u|^2=2\operatorname{Re}\int\overline u\,\frac{\ii}2\Delta u
=-\operatorname{Re}\left(\ii\int|\nabla u|^2\right)=0.
\]
这是空间端的质量守恒证明. $\nabla u$ 也满足同一方程, 所以 $\frac12\int|\nabla u|^2$ 守恒. 对实的、与时间无关的光滑势 $V$, 在衰减充分时 $\ii u_t=-\Delta u/2+Vu$ 的质量及能量 $\int(\tfrac12|\nabla u|^2+V|u|^2)$ 守恒; 分部积分使 Hamilton 算子对称. 时变实势仍保质量, 能量导数为 $\int V_t|u|^2$. 本册不建立一般势的算子定义域和存在理论.

对 $\ii u_t+\tfrac12\Delta u=F$, 初值 $g$, 有
\begin{equation}\label{eq:duhamelschro}
u(t)=U(t)g-\ii\int_0^tU(t-s)F(s)\dd s.
\end{equation}
若 $F\in C([0,T];\Sch)$ 且所需半范数一致受控, 可以求导核验; 若 $F\in L^1([0,T];L^2)$, Bochner 积分给连续 $L^2$ mild 解, 由近似传递弱 PDE. 范数满足 $\|u(t)\|_2\le\|g\|_2+\int_0^t\|F(s)\|_2\dd s$. 指数必须为 $\ee^{-2\pi^2\ii(t-s)|\xi|^2}$, 源项前必须为 $-\ii$, 与热方程相比较时不能沿用正实指数或漏掉 $\ii$.

\originalchapter{变分原理与场的守恒律}\label{ch:lagrangian}
\sourcelecture{21--23}
\originalsection{作用量与 Euler--Lagrange 方程}
\begin{definition}
拉格朗日函数 $\mathcal L(x,u,p)$ 取决于点、标量场与其一阶导数变量 $p=(p_0,\ldots,p_n)$. 在有界时空区域 $K$ 上的作用量是 $A[u]=\int_K\mathcal L(x,u,\partial u)\dd x$. 若对每个 $\psi\in C_c^\infty(K)$ 有 $\frac{\mathrm d}{\mathrm d\eps}A[u+\eps\psi]|_{\eps=0}=0$, 称 $u$ 为驻点.
\end{definition}
驻点是函数而非时空中一个点, 也不必是极小值. 波方程作用量含正负二次项, 不能以正定最小化直觉替代驻点条件.
\begin{theorem}\label{thm:EL}
$\mathcal L\in C^2$, $u\in C^2$ 时, 驻点条件等价于
\[
\partial_\mu\left(\frac{\partial\mathcal L}{\partial p_\mu}(x,u,\partial u)\right)
-\frac{\partial\mathcal L}{\partial u}(x,u,\partial u)=0.
\]
\end{theorem}
\begin{proof}
紧支撑变分使积分内导数连续可积, 链式法则给 $\delta A=\int(\mathcal L_u\psi+\mathcal L_{p_\mu}\partial_\mu\psi)$. 积分分部无边界项, 得 $\delta A=\int[\mathcal L_u-\partial_\mu\mathcal L_{p_\mu}]\psi$. 若连续括号函数非零, 在其同号小邻域取非负光滑测试函数, 积分非零, 矛盾. 反过来括号恒零则每个变分为零.
\end{proof}
取 $\mathcal L=-\tfrac12m^{\mu\nu}u_\mu u_\nu-V(u)=\tfrac12u_t^2-\tfrac12|\nabla u|^2-V(u)$, 则 $\mathcal L_{p_\mu}=-\partial^\mu u$, $\mathcal L_u=-V'(u)$, 所以 $\square u-V'(u)=0$, 等价于 $u_{tt}-\Delta u+V'(u)=0$. $V(u)=m_0^2u^2/2$ 给 Klein--Gordon 方程. 参数 $m_0$ 与度量 $m$ 是不同对象.
\originalsection{坐标变化与度量变分}
光滑向量场 $Y$ 的局部流满足 $\frac{\mathrm d}{\mathrm d\eps}F_\eps(x)=Y(F_\eps(x))$, $F_0(x)=x$. ODE 唯一性给 $F_{\eps+s}=F_\eps\circ F_s$ 及逆 $F_{-\eps}$. 在固定紧集上 Taylor 展开为 $F_\eps=x+\eps Y+O(\eps^2)$, 导数矩阵 $M=I+\eps DY+O(\eps^2)$. 因 $(I+\eps A)^{-1}=I-\eps A+O(\eps^2)$,
\[
M^{-1}=I-\eps DY+O(\eps^2),\qquad \det M^{-1}=1-\eps\diver Y+O(\eps^2).
\]
行列式的一次项是迹: Leibniz 展开中恒等排列贡献 $1+\eps\sum_jA_{jj}$, 其他排列要至少两项非对角因子, 为 $O(\eps^2)$.

新坐标 $\widetilde x=F_\eps(x)$ 下, 标量场 $\widetilde u(\widetilde x)=u(x)$, 协变导数和度量分别变为 $\widetilde p=M^{-T}p$, $\widetilde m=M^{-T}mM^{-1}$. 逆度量 $\widetilde m^{-1}=Mm^{-1}M^T$. 于是 $p^Tm^{-1}p$ 为标量不变量. 坐标变换后度量分量一般不再是固定对角矩阵; 任意微分同胚的不变性要同时变换度量, 与固定 Minkowski 矩阵的 Lorentz 对称性不同.

若度量以参数变化, 对 $mm^{-1}=I$ 求导, 得
\begin{equation}\label{eq:inversevariation}
\delta m^{-1}=-m^{-1}(\delta m)m^{-1}.
\end{equation}
对称度量变分可直接用矩阵等式, 避免独立非对角分量的重复计数.
\originalsection{能量动量张量的变分来源}
对只依赖 $u,\partial u,m$ 且为标量的 $\mathcal L$, 定义 $P^{\mu\nu}$ 为对称矩阵导数, 满足 $\delta\mathcal L=P^{\mu\nu}\delta m_{\mu\nu}$, 在所有有序指标上求和. 令
\[
T^{\mu\nu}=2P^{\mu\nu}+m^{\mu\nu}\mathcal L.
\]
\begin{theorem}\label{thm:stressvariation}
若 $\mathcal L$ 没有显式坐标依赖, 按上述同时变换场与度量的方式坐标不变, 且 $u$ 满足 Euler--Lagrange 方程, 则常 Minkowski 度量下 $\partial_\mu T^{\mu\nu}=0$.
\end{theorem}
\begin{proof}
取紧支撑 $Y$, 新坐标固定时, 标量与导数变化为 $\delta u=-Y^\alpha u_\alpha$, $\delta p_\mu=-\partial_\mu(Y^\alpha u_\alpha)$, 度量变化为 $\delta m_{\mu\nu}=-m_{\alpha\nu}\partial_\mu Y^\alpha-m_{\mu\alpha}\partial_\nu Y^\alpha$. Jacobian 变化为 $-\partial_\alpha Y^\alpha$. 作用量换元后的值不变, 所以
\[
0=\int\left[-\mathcal L_uY^\alpha u_\alpha
-\mathcal L_{p_\mu}\partial_\mu(Y^\alpha u_\alpha)
-2P^{\mu\nu}m_{\alpha\nu}\partial_\mu Y^\alpha
-\mathcal L\partial_\alpha Y^\alpha\right]\dd x.
\]
头两项分部积分后由 Euler--Lagrange 消失, 后两项合为 $-\int T^{\mu\nu}m_{\alpha\nu}\partial_\mu Y^\alpha$. 再分部积分得 $\int(\partial_\mu T^{\mu\nu})m_{\alpha\nu}Y^\alpha=0$. 任意紧支撑 $Y$ 及 $m$ 可逆给散度为零.
\end{proof}
对 $\mathcal L=-Q/2-V(u)$, 由\eqref{eq:inversevariation}得 $P^{\mu\nu}=\tfrac12\partial^\mu u\partial^\nu u$, 因而
\[
T^{\mu\nu}=\partial^\mu u\partial^\nu u-m^{\mu\nu}(Q/2+V(u)),\quad
T^{00}=\tfrac12(u_t^2+|\nabla u|^2)+V(u).
\]
直接求散度也得 $(\square u-V'(u))\partial^\nu u$, 与变分证明相符. 若 $V\ge0$ 则能量密度非负, 但其控制哪些范数仍要看 $V$ 的增长.
\originalsection{对称性与 Noether 流}
若驻点场有一参数内部变换 $u_\eps=u+\eps\eta+O(\eps^2)$, 且拉格朗日变化为全散度 $\delta\mathcal L=\partial_\mu B^\mu$, 则
\[
\partial_\mu(\mathcal L_{p_\mu}\eta-B^\mu)=0.
\]
证明是链式法则 $\delta\mathcal L=\mathcal L_u\eta+\mathcal L_{p_\mu}\partial_\mu\eta$, 用 Euler--Lagrange 代入为 $\partial_\mu(\mathcal L_{p_\mu}\eta)$. 对多个实分量逐个求和即可. 这解释了守恒量来自连续对称性, 并非每个 PDE 都自动有正定能量.
\begin{example}
对波方程加常数的对称性, 求对应流.
\end{example}
\begin{solution}
$V=0$ 时 $\mathcal L$ 只依赖导数, $u\mapsto u+\eps$ 保持它, 所以 $\eta=1$, $B=0$, 流为 $-\partial^\mu u$. 守恒等式就是 $-\square u=0$. 时间平移的流则给二次能量, 是不同对称性.
\end{solution}
\begin{exercise}
$V(u)=m_0^2u^2/2$, 证明足够衰减的 Klein--Gordon 解有 $\int(u_t^2+|\nabla u|^2+m_0^2u^2)$ 守恒.
\end{exercise}
\begin{solution}
将 $u_{tt}-\Delta u+m_0^2u=0$ 乘 $u_t$ 积分. 前后两项分别为 $\tfrac12\partial_t\int u_t^2$ 与 $\tfrac12\partial_t\int m_0^2u^2$, 中间经分部积分为 $\tfrac12\partial_t\int|\nabla u|^2$. 总导数为零. 充分衰减或齐次 Dirichlet/Neumann 边界使边界通量消失. Robin 边界应另计边界能量, 见第九章.
\end{solution}

\originalchapter{输运方程与 Burgers 奇性}\label{ch:transport}
\sourcelecture{1,24; 激波与熵条件为编者补充}
\originalsection{特征流与入流数据}
\begin{theorem}\label{thm:transport}
设 $b(t,x)$ 连续且对空间 $C^1$, 在有限时间区间对空间全局 Lipschitz, 并具有至多线性增长. 对 $u_t+b\cdot\nabla u+a u=F$, 光滑初值 $u(0)=g$ 的经典解沿流 $X'(t)=b(t,X)$ 满足
\[
u(t,X(t;0,y))=\ee^{-\int_0^ta(s,X(s))\dd s}g(y)
+\int_0^t\ee^{-\int_s^ta(r,X(r))\dd r}F(s,X(s))\dd s.
\]
系数与数据足够 $C^1$ 时该式给 $C^1$ 解, 并唯一.
\end{theorem}
\begin{proof}
先修 ODE 存在唯一性和线性增长保证流在给定有限时间区间存在. 导数矩阵满足 $\partial_tD_yX=(D_xb)D_yX$, 初值 $I$; 倒流给逆, 所以流为 $C^1$ 微分同胚. 链式法则给沿曲线的标量 ODE $\frac{\mathrm d}{\mathrm dt}u+a u=F$, 积分因子解即所示公式. 反向由每点唯一的流线确定初始点, 求导核验 PDE. 任意解都服从同一 ODE, 给唯一性.
\end{proof}
有界域的特征若从侧边进入, 要在入流边界 $b\cdot\nu<0$ 给值; 出流边界不能再任意独立指定值. 数据曲面 $\Sigma$ 对时空向量 $(1,b)$ 非特征的条件是该向量不切于 $\Sigma$. 这是第一章直线例子的几何推广.

对齐次输运, $u$ 沿流恒定但体积未必守恒. Jacobian $J=\det D_yX$ 满足 $J'=(\diver b)(t,X)J$, $J(0)=1$, 由行列式微分与线性矩阵 ODE 得到. $\diver b=0$ 时保持体积, 从而保持所有有限 $L^p$ 范数. 连续性方程 $\rho_t+\diver(b\rho)=0$ 则是输运加 $a=\diver b$, 满足 $\rho(t,X)J=\rho(0,y)$, 总质量守恒.
\originalsection{Burgers 的隐式解与破裂时间}
\begin{theorem}\label{thm:burgers}
设 $g\in C^1(\R)$ 有界且 $g'$ 有界. 无粘 Burgers 方程 $u_t+uu_x=0$, $u(0)=g$ 的特征为 $x=y+tg(y)$, 沿特征 $u=g(y)$. 只要 $1+tg'(y)>0$ 且有统一正下界, 该映射是全直线 $C^1$ 微分同胚, 解为 $u(t,x)=g(Y(t,x))$, 且
\[
u_x(t,x)=\frac{g'(Y)}{1+tg'(Y)},\qquad u_t=-u u_x.
\]
若 $m=\inf g'<0$, 则在 $0\le t<T_*=-1/m$ 保持 $C^1$; 若负最小值取得, 则在该特征上 $u_x\to-\infty$ 当 $t\uparrow T_*$. 若 $g'\ge0$, 则对所有正时间存在经典解.
\end{theorem}
\begin{proof}
沿 $X'=u(t,X)$ 有 $\frac{\mathrm d}{\mathrm dt}u(t,X)=u_t+uu_x=0$, 故 $u=g(y)$, $X=y+tg(y)$. 当导数 $1+tg'>0$ 有统一下界, 映射严格递增; $g$ 有界使 $X(y)\to\pm\infty$, 所以满射. 逆函数定理给 $Y_x=(1+tg'(Y))^{-1}$, $Y_t=-g(Y)/(1+tg'(Y))$, 得导数及方程. $t<T_*$ 时下界 $1+tm>0$. 在取得 $m$ 的点分母趋零而分子负非零, 得梯度爆破. 非负导数给下界1, 任意有限时间均可构造. 唯一性在经典且所需流线存在的类别由特征给出.
\end{proof}
更一般的不有界初值需另核全局流与映射满射, 本册不省略该条件. 梯度破裂发生时函数值仍沿特征保持原幅值. 因而守恒一个 $L^2$ 范数并不能防止导数爆破.
\begin{example}
初值 $g(x)=-\tanh x$, 求首次破裂时间.
\end{example}
\begin{solution}
$g'=-\operatorname{sech}^2x$, 最小值 $-1$ 于0取得, 所以 $T_*=1$. 中心特征为 $x=0$, 上面 $u=0$, 但 $u_x(t,0)=-1/(1-t)$ 趋负无穷. $g$ 本身不在 $L^2$, 此例用于梯度奇性, 不援引总 $L^2$ 能量. 可另取具有负斜率的光滑紧支撑初值得到有限能量版本.
\end{solution}
\originalsection{守恒律、激波与熵}
光滑 Burgers 解满足 $u_t+(u^2/2)_x=0$. 乘 $u$ 得 $\partial_t(u^2/2)+\partial_x(u^3/3)=0$. 对紧支撑或足够可积且允许求导和消失边界通量的解, $\int u^2$ 守恒. 仅“在无穷远趋零”不保证 $L^2$ 可积.

破裂后采用分布守恒律. $u$ 为局部有界弱解, 若对 $\phi\in C_c^\infty((0,\infty)\times\R)$ 有 $\int[u\phi_t+(u^2/2)\phi_x]=0$. 初值以局部 $L^1$ 迹给定.
\begin{proposition}\label{prop:RH}
分段常值 $u=u_L$ 于 $x<st$, $u=u_R$ 于 $x>st$, 是 Burgers 弱解当且仅当 $s=(u_L+u_R)/2$, 若 $u_L\ne u_R$.
\end{proposition}
\begin{proof}
对移动界面两侧分部积分, 界面分布项为 $[-s(u_R-u_L)+(u_R^2-u_L^2)/2]\delta_{x=st}$. 系数为零即 Rankine--Hugoniot 条件 $s[u]=[u^2/2]$, 除以非零跳跃得速度. 常值两侧的内部导数为零.
\end{proof}
弱守恒不能单独确定物理解. 对凸熵 $\eta$, 熵通量满足 $q'(u)=\eta'(u)u$, 熵条件要求 $\partial_t\eta(u)+\partial_xq(u)\le0$ 于分布意义. 对平方熵 $\eta=u^2/2$, $q=u^3/3$, 激波系数为
\[
-s[\eta]+[q]=\frac{(u_R-u_L)^3}{12}.
\]
所以下降跃迁 $u_L>u_R$ 满足非正熵产生, 上升跃迁不满足. 一般凸熵也可直接核验. 置 $a=u_L>b=u_R$, $s=(a+b)/2$, 其界面系数为
\[
-s[\eta]+[q]=-\int_b^a\eta'(v)(v-s)\dd v
=-\int_b^a[\eta'(v)-\eta'(s)](v-s)\dd v\le0.
\]
第二等号用 $\int_b^a(v-s)\dd v=0$, 最后一步用凸性使 $\eta'$ 单调. 因而下降激波满足所有光滑凸熵条件; 一般凸熵可由光滑凸函数逼近. 两侧特征均向界面汇聚.
\begin{example}
初值在零点左侧为 $u_L$, 右侧为 $u_R$, 求熵解.
\end{example}
\begin{solution}
若 $u_L>u_R$, 取速度 $s=(u_L+u_R)/2$ 的下降激波. 若 $u_L<u_R$, 取稀疏波
\[
u(t,x)=\begin{cases}u_L,&x/t\le u_L,\\x/t,&u_L<x/t<u_R,\\u_R,&x/t\ge u_R.\end{cases}
\]
扇形内部 $u_t=-x/t^2$, $uu_x=x/t^2$, 互相抵消; 两连接线函数值连续, 不产生跳跃分布, 因而为弱解. 它在每个 $t>0$ 局部 Lipschitz, 链式法则给熵等式几乎处处, 连续连接使熵无奇异贡献, 故满足凸熵条件. 初时在局部 $L^1$ 中趋原阶梯函数. 本册证明这些基本解, 不建立任意有界数据熵解的完整存在唯一性理论.
\end{solution}

\originalchapter{课程作业: 基础与热方程}
\label{ch:assignheat}
本章至课程考试部分均保留 MIT 18.152, Fall 2011, Jared Speck 的原题编号. 公开完整题面译为中文; 只有期中附有官方答案, 其余解答为 AI 独立编写、经数学审读的解答. 为使积分和解的意义明确, 增补条件及原卷勘误均明示. 题目链接与缺项表见末章.
\originalsection{作业1: 2011年9月15日交}
\begin{exercise}\label{ps:1I}
\textbf{原题 I.} 在有界光滑域 $\Omega\subset\R^n$, $u,v\in C^2(\overline\Omega)$, 证明 Green 第二恒等式. 原题未写有界性; 无界域需另有积分与无穷远通量条件.
\end{exercise}
\begin{solution}
令 $F=u\nabla v-v\nabla u$. 乘积法则给 $\diver F=u\Delta v-v\Delta u$, 交叉项 $\nabla u\cdot\nabla v$ 抵消. 散度定理及 $F\cdot\nu=u\partial_\nu v-v\partial_\nu u$ 给
\[
\int_\Omega(u\Delta v-v\Delta u)\dd x
=\int_{\partial\Omega}(u\partial_\nu v-v\partial_\nu u)\dd S.
\]
\end{solution}
\begin{exercise}\label{ps:1II}
\textbf{原题 II.} 对 $0<\eps<1/2$, 证明 $h(x)=\sin x\log(1+x^2)/|x|^{1-\eps}$ 属于 $L^2(\R)$.
\end{exercise}
\begin{solution}
在 $|x|\le1$ 用 $|\sin x|\le|x|$, $\log(1+x^2)\le x^2$, 得 $|h|^2\le|x|^{4+2\eps}$ 可积. 在 $|x|\ge1$, $|h|^2\le\log^2(1+x^2)/|x|^{2-2\eps}$. 取 $0<\delta<1-2\eps$, 则 $\log^2(1+x^2)\le C_\delta |x|^\delta$; 指数 $2-2\eps-\delta>1$, 两个尾部均可积. 在0定义为0不改变 $L^2$ 类.
\end{solution}
\begin{exercise}\label{ps:1III}
\textbf{原题 III.} 实内积空间 $V$ 中证明 Cauchy--Schwarz 与三角不等式.
\end{exercise}
\begin{solution}
$w=0$ 显然. 否则 $q(t)=\|v+tw\|^2=\|v\|^2+2t\langle v,w\rangle+t^2\|w\|^2\ge0$. 于 $t=-\langle v,w\rangle/\|w\|^2$ 取最小值, 得 $\langle v,w\rangle^2\le\|v\|^2\|w\|^2$. 再展开 $\|v+w\|^2$ 并用刚得不等式, 得 $\|v+w\|^2\le(\|v\|+\|w\|)^2$, 开平方即可.
\end{solution}
\begin{exercise}\label{ps:1IV}
\textbf{原题 IV.} 对实值 $L^2(\R)$ 函数的几乎处处等价类, 验证 $\langle f,g\rangle=\int fg$ 是内积, 推出积分 Cauchy--Schwarz. 并证明
\[
\left|\int_\R\frac{\sin x\log(1+x^2)}{|x|^{3/4}}f(x)\dd x\right|\le C\|f\|_2.
\]
\end{exercise}
\begin{solution}
$2|fg|\le f^2+g^2$ 先保证积分存在. 线性和对称性由积分性质得到. $\int f^2\ge0$, 等号当且仅当 $f=0$ 几乎处处, 即零等价类. 应用原题 III 得 $|\int fg|\le\|f\|_2\|g\|_2$. 原题 II 取 $\eps=1/4$, 令上述权函数为 $h$, 得 $C=\|h\|_2<\infty$. 复值版内积应改为 $\int f\overline g$.
\end{solution}
原题 V 是阅读 Salsa 的附录 A, 不是公开完整习题, 在末章登记.
\originalsection{作业2: 2011年9月22日交}
原题 I、II 仅引用 Salsa 题号, 未收入猜造题面.
\begin{exercise}\label{ps:2III}
\textbf{原题 III.} 对零 Dirichlet 热流
\[
u(t,x)=\sum_{m=1}^\infty\frac{2(-1)^{m+1}}{m\pi}\ee^{-m^2\pi^2t}\sin(m\pi x),
\]
证明 $\|u(t,\cdot)-x\|_{L^2(0,1)}\to0$ 当 $t\downarrow0$.
\end{exercise}
\begin{solution}
$x$ 的正弦系数为 $b_m=2\int_0^1x\sin(m\pi x)\dd x=2(-1)^{m+1}/(m\pi)$. 正弦系完备及 Parseval 给
\[
\|u(t)-x\|_2^2=\tfrac12\sum_{m\ge1}|b_m|^2|\ee^{-m^2\pi^2t}-1|^2.
\]
每项趋0, 并受可和的 $|b_m|^2/2$ 控制, 所以趋0. 初值在 $x=1$ 为1, 正时间边值为0, 角点不连续, 因此不能把此解称为闭矩形连续解. $L^2$ 初始迹仍成立.
\end{solution}
\begin{exercise}\label{ps:2IV}
\textbf{原题 IV.} $u_t=u_{xx}$ 于 $(0,\infty)\times(0,\ell)$, 零 Dirichlet 边值, 初值 $\ell^{-2}x(\ell-x)$. 不显式求解, 依次求初始 $L^2$ 范数、能量导数、$C^0$ 与 Poincaré 型估计, 证明 $\|u(t)\|_2\le\sqrt{\ell/30}\ee^{-t/\ell^2}$.
\end{exercise}
\begin{solution}
代 $x=\ell y$ 得 $\|u(0)\|_2^2=\ell\int_0^1y^2(1-y)^2\dd y=\ell/30$. 正时间分部积分给 $\partial_t\|u\|_2^2=2\int u u_{xx}=-2\|u_x\|_2^2$. 对 $u(0)=0$ 的空间端点条件,
$|u(x)|=|\int_0^xu_x|\le\sqrt\ell\|u_x\|_2$, 所以 $\|u\|_\infty\le\sqrt\ell\|u_x\|_2$, 再积分得 $\|u\|_2^2\le\ell^2\|u_x\|_2^2$. 故 $y(t)=\|u(t)\|_2^2$ 满足 $y'\le-2\ell^{-2}y$. 从 $\delta>0$ 积分后令 $\delta\downarrow0$, 用 $L^2$ 初始连续性得 $y(t)\le(\ell/30)\ee^{-2t/\ell^2}$.

原卷末段的区间误写 $[0,1]$, 此处统一为 $(0,\ell)$. 初始角点高阶相容性不满足, 所用类别为正时间光滑且具有初值迹, 不强求 $C^{1,2}$ 延伸到角点.
\end{solution}
\begin{exercise}\label{ps:2V}
\textbf{原题 V.} $D>0$, 从 $u(t,x)=(Dt)^{-1/2}V(x/\sqrt{Dt})$ 推导非负、偶、质量1且两端衰减的一维热基本解.
\end{exercise}
\begin{solution}
置 $\zeta=x/\sqrt{Dt}$. 求导代入 $u_t-Du_{xx}=0$, 得 $V''+(\zeta V)'/2=0$, 即 $(V'+\zeta V/2)'=0$. 偶性给 $V'(0)=0$, 故积分常数为0. 分离变量得 $V=V(0)\ee^{-\zeta^2/4}$. 归一化 $1=\int V\dd\zeta=\sqrt{4\pi}V(0)$, 所以 $u=(4\pi Dt)^{-1/2}\ee^{-x^2/(4Dt)}$. 此式确非负、偶、衰减, 且直接求导验证 PDE.
\end{solution}

\originalsection{作业3: 2011年9月29日交}
\begin{exercise}\label{ps:3I}
\textbf{原题 I.} 零边值热流的初值为 $f(x)=x(1-x)$, 证明 $u(t,x)=u(t,1-x)$.
\end{exercise}
\begin{solution}
$v(t,x)=u(t,1-x)$ 满足 $v_t=v_{xx}$, 两端为0, 初值 $f(1-x)=f(x)$. 对差 $w=u-v$ 用热流能量唯一性或最大原理得 $w=0$. 类别为正时间光滑且闭域连续, 不强求不相容的角点二阶导数.
\end{solution}
\begin{exercise}\label{ps:3II}
\textbf{原题 II, 含勘误.} 初值仍为 $f=x(1-x)$, 两端边值均为常数, $u$ 在闭域连续. 证明 $u\ge0$, 求所有 $\alpha,\beta>0$ 使 $u\le\alpha f\ee^{-\beta t}$, 并说明长期一致趋零. 原卷最后的 $t\to0$ 应为 $t\to\infty$.
\end{exercise}
\begin{solution}
闭域连续性及 $f(0)=f(1)=0$ 强制两个边界常数为0. 最大原理给非负. 一个简明充分范围是 $\alpha\ge1$, $0<\beta\le8$: 令 $w=\alpha f\ee^{-\beta t}$, 则 $w_t-w_{xx}=\alpha\ee^{-\beta t}(2-\beta f)\ge0$, 因 $f\le1/4$; 初边值比较给所需界.

题面问“所有”, 所以上述范围不足. 给出以下精确的谱式刻画, 其中只含显式数列而没有未知 PDE 解:
\[
A(\beta)=\sup_{t\ge0,\ 0<x<1}
\frac{8\sum_{j=0}^\infty(2j+1)^{-3}\ee^{-((2j+1)^2\pi^2-\beta)t}\sin((2j+1)\pi x)}{\pi^3x(1-x)}.
\]
全部答案恰为 $0<\beta\le\pi^2$, $\alpha\ge A(\beta)$. 此隐式常数由显式收敛级数定义, 适用于没有简洁初等闭式的参数部分. 下面证明其有限性、必要性与充分性.

两次分部积分给 $f$ 的正弦系数 $8/(m^3\pi^3)$ 于奇数 $m$, 偶数为0. 热级数因此正是上式乘 $f\ee^{-\beta t}$. 对 $t>0$ 绝对收敛, 对 $t=0$ 以连续初值解释. 由于 $f\le\tfrac12\sin\pi x$, 热比较给 $u\le\tfrac12\ee^{-\pi^2t}\sin\pi x$. 此初始不等式可用 $\sin\pi x\ge2x(1-x)$ 证明: 在 $[0,1/2]$ 凹性给 $\sin\pi x\ge2x\ge2x(1-x)$, 另半区对称. 又 $\sin\pi x\le\pi\min(x,1-x)$, 故 $\sin\pi x/f\le2\pi$. 对 $\beta\le\pi^2$ 得 $1\le A(\beta)\le\pi$; 下界由初时比值为1. 所以定义确给有限阈值, 逐点取上确界证明 $\alpha\ge A$ 的充要性.

若 $\beta>\pi^2$, 把假设界乘正的 $\sin\pi x$ 积分. 左侧第一系数为 $\ee^{-\pi^2t}\int f\sin\pi x$, 右侧为 $\alpha\ee^{-\beta t}\int f\sin\pi x$, 矛盾于大时间. 当 $\beta\le8$, 上面超解证明与初值下界给 $A=1$. 当 $\beta>8$, 在中心 $u_t(0,1/2)=f''=-2$, 故 $\partial_t(\ee^{\beta t}u/f)|_{t=0}=\beta-8>0$, 得 $A>1$. 例如 $\beta=\pi^2,\alpha=\pi$ 可行, 说明仅报 $\beta\le8$ 会漏解. 最后任取可行参数, $\|u(t)\|_\infty\le\alpha\ee^{-\beta t}/4\to0$.
\end{solution}
原题 III、IV 只有 Salsa 引用, 见缺项表.
\begin{exercise}\label{ps:3V}
\textbf{原题 V.} $L=\partial_t+\widetilde L$, 其中 $\widetilde L$ 线性且仅作用于空间. 若 $v^{(s)}(\tau,x)$ 满足 $Lv^{(s)}=0$, $v^{(s)}(0,x)=F(s,x)$, 证明 $v(t,x)=\int_0^tv^{(s)}(t-s,x)\dd s$ 满足 $Lv=F$, $v(0)=0$.
\end{exercise}
\begin{solution}
假设参数族足够光滑且允许积分求导. Leibniz 公式给 $v_t=v^{(t)}(0,x)+\int_0^t\partial_\tau v^{(s)}(t-s,x)\dd s$. 线性及空间算子换积分给 $Lv=F(t,x)+\int_0^t(\partial_\tau+\widetilde L)v^{(s)}(t-s,x)\dd s=F(t,x)$. 零长度积分给初值0. 原卷齐次初值写成 $f(t,x)$ 留有自由时间, 应改为一般 $h(x)$; 参数族的初值才是 $F(s,x)$.
\end{solution}
\begin{exercise}\label{ps:3VI}
\textbf{原题 VI.} 给热方程 $u_t-D\Delta u=F$, $u(0)=g$ 的表示公式, 并选充分条件证明公式有意义且正确.
\end{exercise}
\begin{solution}
取 $g\in C_c^\infty(\R^n)$, $F\in C^\infty([0,T]\times\R^n)$ 且空间支撑于同一紧集. 定义
\[
u(t,x)=\Gamma_D(t)*g(x)+\int_0^t\Gamma_D(t-s)*F(s)(x)\dd s.
\]
第一项满足齐次 PDE 并以近似恒等趋 $g$. 源项可将空间导数移到 $F$: $\Delta(\Gamma*F)=\Gamma*\Delta F$, 后者一致有界, 避免核在 $s=t$ 的导数奇性. 对时间求导得到端点 $F(t,x)$ 与 $D\Delta$ 积分项, 即非齐次方程. 源项的 $L^\infty$ 范数不超过 $t\sup_s\|F(s)\|_\infty$, 初时趋0. 齐次项与源项相加给初值和 PDE, 有界类唯一性由热最大原理给出.
\end{solution}

\originalchapter{课程作业: 势论与波能量}
\originalsection{作业4: 2011年10月6日交}
原题 I--V 是商业教材题号引用, 见缺项表.
\begin{exercise}\label{ps:4VI}
\textbf{原题 VI.} 单位球中 $\Delta u=f$, $u|_{\partial B}=g$, 证明存在与 $f,g$ 无关的 $C$ 使 $\|u\|_\infty\le C(\|f\|_\infty+\|g\|_\infty)$. 补充 $f$ 有界、$u$ 有连续边界迹, 以 $\sup$ 解释范数.
\end{exercise}
\begin{solution}
置 $M=\|g\|_\infty$, $F=\|f\|_\infty$, $b=(1-|x|^2)/(2n)$, 则 $b\ge0$, $\Delta b=-1$. 对 $v=u-M-Fb$ 有 $\Delta v=f+F\ge0$, 边界 $v\le0$, 最大原理给 $u\le M+Fb$. 同样用于 $-u$ 得 $|u|\le M+F/(2n)$. 可取 $C=\max(1,1/(2n))=1$. 开球光滑并不自动使 $f$ 有界, 此条件不可省.
\end{solution}
\originalsection{作业5: 2011年10月13日交}
本节沿原作业的符号 $\Delta_yG(x,y)=\delta_x$, 故 $G$ 为负; 与正文的 $-\Delta$ 正 Green 核相反.
\begin{exercise}\label{ps:5I}
\textbf{原题 I.} 三维单位球的核
\[
G(x,y)=-\frac1{4\pi|x-y|}+\frac1{4\pi|x||x/|x|^2-y|},\quad x\ne0,
\]
以及 $G(0,y)=-1/(4\pi|y|)+1/(4\pi)$, 直接证明 $G\le0$.
\end{exercise}
\begin{solution}
计算 $|x|^2|x/|x|^2-y|^2-|x-y|^2=(1-|x|^2)(1-|y|^2)>0$. 所以像点项的分母较大, 正项小于负项的绝对值, $G<0$ 于 $x\ne y$. $x=0$ 时 $|y|<1$ 直接得到负性. 对角线上核发散为 $-\infty$, 不作有限点值断言.
\end{solution}
\begin{exercise}\label{ps:5II}
\textbf{原题 II.} 计算 $\int_B G(x,y)\dd y$, 推出 $\int_B|G|\le1/6$, 并从表示式证明 $\Delta u=f$, 边值 $g$ 的一致数据估计.
\end{exercise}
\begin{solution}
$w(x)=(|x|^2-1)/6$ 满足 $\Delta w=1$, 边值0. Green 表示式给 $w(x)=\int_BG(x,y)\dd y$. 由负性,
$\int_B|G|=(1-|x|^2)/6\le1/6$. 又 $P(x,\sigma)=\partial_{\nu_\sigma}G=(1-|x|^2)/(4\pi|x-\sigma|^3)\ge0$. 将常函数1代入表示式得 $\int_{\partial B}P=1$. 因此
$|u(x)|\le\|f\|_\infty/6+\|g\|_\infty$. 仍须 $f$ 有界且边界迹连续; 不从开球光滑推有限最大值.
\end{solution}
\begin{exercise}\label{ps:5III}
\textbf{原题 III.} $u$ 在 $\R^3$ 调和, $|u(x)|\le\log(1+|x|)$, 证明 $u\equiv0$.
\end{exercise}
\begin{solution}
$u(0)=0$. 在 $B_R$ 中 $v_R=u+\log(1+R)\ge0$ 且调和. 对固定 $x$, $r=|x|<R$, 球 Harnack 给
\[
\frac{R(R-r)}{(R+r)^2}v_R(0)\le v_R(x)\le
\frac{R(R+r)}{(R-r)^2}v_R(0).
\]
两个系数均为 $1+O(r/R)$, 故 $|u(x)|\le C_x\log(1+R)/R$ 当 $R$ 足够大. 令 $R\to\infty$ 得0. 原题 IV、V 的完整商业书题面未公开, 本册不依据提示猜题.
\end{solution}
\originalsection{作业6: 2011年10月20日交}
\begin{exercise}\label{ps:6I}
\textbf{原题 I(a,b).} 对 $u_{tt}=\Delta u$, 置 $J^0=(u_t^2+|\nabla u|^2)/2$, $J^i=-u_tu_i$. 证明 $\partial_\mu J^\mu=0$, 且对欧氏点积、$V=(1,\omega)$, $|\omega|\le1$, 有 $V\cdot J\ge0$.
\end{exercise}
\begin{solution}
(a) 求导后 $\partial_tJ^0+\partial_iJ^i=u_t(u_{tt}-\Delta u)=0$, 混合导数项抵消.
(b) $V\cdot J=(u_t^2+|\nabla u|^2)/2-u_t\omega\cdot\nabla u\ge(|u_t|-|\nabla u|)^2/2\ge0$, 用 $|\omega\cdot\nabla u|\le|\nabla u|$.
\end{solution}
\begin{exercise}\label{ps:6II}
\textbf{原题 II(a--d).} 在后向截锥 $0\le\tau\le t\le R$, $|y-p|\le R-\tau$, 求侧面外单位法向, 证明局部能量 $\int_{B_{R-t}(p)}J^0(t)\le\int_{B_R(p)}J^0(0)$; 推出紧支撑数据的有限传播, 再补论证全空间能量等号.
\end{exercise}
\begin{solution}
(a) 侧面是 $h=\tau+|y-p|-R=0$, 内部 $h\le0$, $\nabla_{\tau,y}h=(1,\omega)$, $\omega=(y-p)/|y-p|$. 因长度为 $\sqrt2$, 外法向 $N=(1,\omega)/\sqrt2$.
(b) 时空散度定理给 $0=E_B(t)-E_B(0)+\int_{\rm side}N\cdot J$. 后项非负, 即局部不等式.
(c) 对 $|x-p|>R_0+t$, 选足够小 $\delta>0$ 使 $B_{t+\delta}(x)$ 与初值支撑不交. 局部估计给在 $B_\delta(x)$ 的 $u_t,\nabla u$ 均为0. 亦可将锥缩至任意 $(s,x)$, $0\le s\le t$, 得 $u_t(s,x)=0$. 因 $u(0,x)=0$, 沿竖直线积分得 $u(t,x)=0$. 这一步排除“梯度零但值为非零常数”的漏洞.
(d) 固定有限时间 $T$, 支撑在 $B_{R_0+T}(p)$. 选固定更大的球, 侧面通量为0, 散度定理给 $\int J^0(t)=\int J^0(0)$. 故 $\|\nabla_{t,x}u(t)\|_2=\bigl(\int(|g|^2+|\nabla f|^2)\bigr)^{1/2}$. 仅令局部不等式的 $R\to\infty$ 只得“$\le$”, 必须另用零通量或时间反演补等号.
\end{solution}
\begin{exercise}\label{ps:6III}
\textbf{原题 III.} 一维紧支撑光滑初值于 $[-R,R]$, 证明当 $t\ge R$ 时 $P^2=\int u_x^2$, $K^2=\int u_t^2$ 均等于 $(P^2+K^2)/2$.
\end{exercise}
\begin{solution}
由 d'Alembert 公式,
$u_x=\tfrac12[f'(x+t)+f'(x-t)+g(x+t)-g(x-t)]$,
$u_t=\tfrac12[f'(x+t)-f'(x-t)+g(x+t)+g(x-t)]$.
置 $A=f'(x+t)+g(x+t)$, $B=f'(x-t)-g(x-t)$, 则 $u_x=(A+B)/2$, $u_t=(A-B)/2$, 所以 $P^2-K^2=\int AB$. $A$ 支撑于 $[-R-t,R-t]$, $B$ 于 $[-R+t,R+t]$, $t\ge R$ 时除单点外不交, 积分为0. 总能量由前题守恒.
\end{solution}

\originalchapter{课程作业: Lorentz 几何与共形能量}
\originalsection{作业7: 2011年11月3日交}
\begin{exercise}\label{ps:7I}
\textbf{原题 I.} 三维波方程的光滑初值 $f,g$ 支撑包含于闭球 $\overline B_R$, 证明 $|u(t,x)|\le C/(1+t)$, $C$ 仅依赖 $R,\|f\|_\infty,\|\nabla f\|_\infty,\|g\|_\infty$.
\end{exercise}
\begin{solution}
Kirchhoff 式写成
\[
u=\frac1{4\pi}\int_{\Sph^2}[f(x+t\omega)+t\nabla f(x+t\omega)\cdot\omega+tg(x+t\omega)]\dd\omega.
\]
$t\le2R$ 时取全球面积给 $|u|\le\|f\|_\infty+2R(\|\nabla f\|_\infty+\|g\|_\infty)$, 将常数乘 $1+2R$ 即满足所需界.

$t\ge2R$ 且移动球面与 $B_R$ 相交, 令 $d=|x|$, 则 $t-R\le d\le t+R$, 故 $d\ge t/2$. 球冠面积为 $2\pi t^2(1-\cos\theta_0)$, 其中 $\cos\theta_0=(d^2+t^2-R^2)/(2dt)$. 所以面积 $=\pi t[R^2-(d-t)^2]/d\le2\pi R^2$. 不交时面积0. 由 $\dd\omega=t^{-2}\dd S$, 得 $|u|\le (R^2/2t^2)\|f\|_\infty+(R^2/2t)(\|\nabla f\|_\infty+\|g\|_\infty)$. 在 $t\ge2R$ 把它改写为 $C/(1+t)$ 即可, $R>0$ 固定.
\end{solution}
\begin{exercise}\label{ps:7II}
\textbf{原题 II(a--c).} $\Lambda^Tm\Lambda=m$, $\Omega^Tm\Omega=m$, $m=\operatorname{diag}(-1,1,\ldots,1)$. 证明 $|\det\Lambda|=1$, 乘积与逆仍为 Lorentz 变换.
\end{exercise}
\begin{solution}
(a) 取行列式得 $(\det\Lambda)^2\det m=\det m$, 因 $\det m=-1\ne0$, 结论成立并保证可逆.
(b) $(\Lambda\Omega)^Tm(\Lambda\Omega)=\Omega^T(\Lambda^Tm\Lambda)\Omega=m$.
(c) 原等式两侧左乘 $(\Lambda^{-1})^T$、右乘 $\Lambda^{-1}$ 得 $m=(\Lambda^{-1})^Tm\Lambda^{-1}$.
\end{solution}
\begin{exercise}\label{ps:7III}
\textbf{原题 III.} $|v|<1$, $\gamma=(1-v^2)^{-1/2}$, 证明矩阵 $B$ 的 $(t,x^1)$ 块为 $\left(\begin{smallmatrix}\gamma&-\gamma v\\-\gamma v&\gamma\end{smallmatrix}\right)$、其它空间方向为恒等的 boost 为 Lorentz 变换.
\end{exercise}
\begin{solution}
二阶块计算给 $-(\gamma t-\gamma vx)^2+(\gamma x-\gamma vt)^2=\gamma^2(1-v^2)(-t^2+x^2)=-t^2+x^2$. 其它平方不变. 二次型等式极化后给 $B^TmB=m$. 二阶块行列式 $\gamma^2(1-v^2)=1$, 故总行列式1.
\end{solution}
\begin{exercise}\label{ps:7IV}
\textbf{原题 IV.} 未来类时向量 $X$ 可由行列式1的 Lorentz 变换变为 $(c,0,\ldots,0)$, $c>0$.
\end{exercise}
\begin{solution}
$n\ge2$ 时空间正旋转把 $\vec X$ 变为 $(a,0,\ldots)$, $a=|\vec X|$. $n=1$ 无须旋转, 用带符号的 $a=X^1$. 置 $v=a/X^0$, 类时性给 $|v|<1$. 上题 boost 后空间首分量 $\gamma(a-vX^0)=0$, 时间分量 $\gamma(X^0-va)=\sqrt{(X^0)^2-a^2}>0$. 旋转与 boost 均行列式1, 乘积是所需变换. $\vec X=0$ 时取恒等.
\end{solution}
\begin{exercise}\label{ps:7V}
\textbf{原题 V.} $n\ge1$, 对两个未来类时向量 $X,Y$, 找共同未来零向量 $L,\overline L$, $m(L,\overline L)=-2$, 及正数 $a,b,c,d$ 使 $X=aL+b\overline L$, $Y=cL+d\overline L$.
\end{exercise}
\begin{solution}
用原题 IV 使 $X=(X^0,0)$, 再空间旋转使 $Y=(Y^0,Y^1,0,\ldots)$; $n=1$ 直接已是此形. 取 $L=(1,1,0,\ldots)$, $\overline L=(1,-1,0,\ldots)$, 则零性及互积 $-2$ 直接成立. $a=b=X^0/2>0$, $c=(Y^0+Y^1)/2>0$, $d=(Y^0-Y^1)/2>0$, 因 $Y^0>|Y^1|$. 逆变换回原系; 所用变换保持未来方向、内积和系数, 即结论.
\end{solution}
\originalsection{作业8: 2011年11月10日交}
\begin{exercise}\label{ps:8I}
\textbf{原题 I.} $T_{\mu\nu}=\phi_\mu\phi_\nu-m_{\mu\nu}Q/2$, $Q=m^{\alpha\beta}\phi_\alpha\phi_\beta$. 若时空梯度非零, 证明对未来类时 $X,Y$ 有 $T(X,Y)>0$.
\end{exercise}
\begin{solution}
补全归一零向量对为 $L,\overline L,e_1,\ldots,e_{n-1}$, 其中横向向量正交单位. 设 $A=L\phi$, $B=\overline L\phi$, $H=\sum(e_i\phi)^2$, 则 $Q=-AB+H$. 因此 $T(L,L)=A^2$, $T(\overline L,\overline L)=B^2$, $T(L,\overline L)=AB+Q=H$. 三者非负, 且若全为0, 则基底所有导数0, 与梯度非零矛盾. 按作业7-V写 $X=aL+b\overline L$, $Y=cL+d\overline L$, 所有系数正,
$T(X,Y)=acA^2+bdB^2+(ad+bc)H>0$. 两者过去指向时同时取负, 双线性使结论不变.
\end{solution}
\begin{exercise}\label{ps:8II}
\textbf{原题 II(a--f), 勘误题.} 三维空间中 $K=(1+t^2+r^2,2tx)$, $J_K^\mu=-T^{\mu\nu}K_\nu$. (a)证未来类时; (b)证 $\partial_\mu K_\nu+\partial_\nu K_\mu=4tm_{\mu\nu}$; (c)原题声称 $\operatorname{tr}_mT=0$; (d)原题声称 $\partial_\mu J_K^\mu=0$; (e)求零方向表达; (f)原题声称 $\int J_K^0$ 守恒. 原卷(d,f)还把波方程印作 $m^{\mu\nu}\phi_\mu\phi_\nu=0$, 此处按题意恢复 $\Box\phi=0$. 其中(c,d,f)不成立, 要说明错误并给修正, 不能证明假命题.
\end{exercise}
\begin{solution}
(a) $K^0>0$, 且 $(K^0)^2-|\vec K|^2=[1+(t+r)^2][1+(t-r)^2]>0$.
(b) $K_0=-(1+t^2+r^2)$, $K_i=2tx_i$. 因此 $2\partial_0K_0=-4t$, $\partial_0K_i+\partial_iK_0=2x_i-2x_i=0$, $\partial_iK_j+\partial_jK_i=4t\delta_{ij}$, 给所需式.
(c) 在4维时空 $m_{\mu\nu}T^{\mu\nu}=Q-4Q/2=-Q$, 一般非零. 例如波解 $\phi=t$ 有 $Q=-1$、迹1.
(d) 场方程给 $\partial_\mu T^{\mu\nu}=0$. 用对称性与(b), 得 $\partial_\mu J_K^\mu=-2t\operatorname{tr}_mT=2tQ$, 一般不为0.
(e) $r>0$ 取 $L=\partial_t+\partial_r$, $\overline L=\partial_t-\partial_r$, $H=|\nabla\phi|^2-\phi_r^2$. 置 $A=1+(t+r)^2$, $B=1+(t-r)^2$, 则 $K=(AL+B\overline L)/2$, $\partial_t=(L+\overline L)/2$. 原题 I 的零标架计算给
\[
J_K^0=\tfrac14\{A(L\phi)^2+B(\overline L\phi)^2+2(1+t^2+r^2)H\}.
\]
$r=0$ 以原笛卡尔表达解释, 不赋予任意径向值.
(f) 紧支撑波解满足 $\frac{\mathrm d}{\mathrm dt}\int J_K^0=2t\int Q$, 不是恒0. 例如非零紧支撑 $f$、$g=0$, 初时 $\int Q=\int|\nabla f|^2>0$, 连续性给小正时刻导数非零. 正确修正流为 $\widetilde J^\mu=J_K^\mu-2t\phi\partial^\mu\phi+\phi^2\partial^\mu t$; 其守恒及正平方表达在下节官方 Bonus 全部证明.
\end{solution}

\originalsection{官方 Bonus: 2011年11月29日交}
\begin{exercise}\label{ps:bonus}
\textbf{原题 I(a--i), 原卷两个f依次记为 $f_1,f_2$.} 对作业8的 $K,T,J_K$, 令
\[
\widetilde J^\mu=J_K^\mu-2t\phi\partial^\mu\phi+\phi^2\partial^\mu t.
\]
(a,b)证 $K$ 未来类时及共形 Killing 式; (c)求应力迹; (d)求 $J_K$ 散度; (e)证 $\widetilde J$ 无散度; ($f_1$)求 $J_K^0$ 零标架式; ($f_2$)求 $\widetilde J^0$; (g)证紧支撑 $\phi$ 的径向分部积分式; (h)给正平方表达; (i)证修正共形能量守恒.
\end{exercise}
\begin{solution}
(a,b,$f_1$)与作业8-II(a,b,e)完全重复, 证明见习题\ref{ps:8II}; 保留这里的课程位置而复用相同推导.
(c)4维迹为 $-Q$.
(d)对称性给 $\partial_\mu J_K^\mu=2tQ$.
(e) 在 $\Box\phi=0$ 上,
\[
\partial_\mu(-2t\phi\partial^\mu\phi)=-2\phi\partial^0\phi-2tQ,
\qquad \partial_\mu(\phi^2\partial^\mu t)=2\phi\partial^0\phi.
\]
$\partial^0=-\partial_t$, 交叉项相消, 再加(d)得0.
($f_2$) $\partial^0t=-1$, 所以 $\widetilde J^0=J_K^0+2t\phi\phi_t-\phi^2$.
(g) $\diver_x(x\phi^2)=3\phi^2+2r\phi\phi_r$. 紧支撑使散度积分0, 故 $-\int\phi^2=(2/3)\int r\phi\phi_r$. 也可用球坐标, $\int_0^\infty r^3\partial_r(\phi^2)\dd r=-3\int_0^\infty r^2\phi^2\dd r$.
(h) 置
\[
P=\tfrac14\{(L\phi)^2+[L((t+r)\phi)]^2+(\overline L\phi)^2
+[\overline L((t-r)\phi)]^2+2(1+t^2+r^2)H\}.
\]
这里 $L(t+r)=2$, $\overline L(t-r)=2$. 展开两个平方得 $P=J_K^0+2t\phi\phi_t+2r\phi\phi_r+2\phi^2$, 因而
\[
P-\widetilde J^0=2r\phi\phi_r+3\phi^2=\diver_x(x\phi^2),\qquad
\int P\dd x=\int\widetilde J^0\dd x\ge0.
\]
原卷(h)写为点态相等, 实际只有积分后相等, 或相差上述空间散度. 不把分部积分后的等式误称逐点等式.
(i) 若光滑初值紧支撑, 波方程有限传播保证每个有限时间段的统一紧支撑. 时空柱的侧面通量0, 无散度给 $\int\widetilde J^0(t)=\int\widetilde J^0(0)$. 用(h)即
\[
\mathcal E_{\rm conf}(t)=\int_{\R^3}P(t,x)\dd x=\int_{\R^3}P(0,x)\dd x.
\]
右端以 $f=\phi(0)$、$g=\phi_t(0)$ 代入 $L\phi=g+f_r$, $\overline L\phi=g-f_r$, $H=|\nabla f|^2-f_r^2$ 即由初值完全确定. 一般非紧支撑场须加有限加权能量及无穷远通量条件. 本节建立守恒, 不据此声称已证明所有零导数的精细衰减率.
\end{solution}

\originalchapter{课程作业: Fourier 方法与场论}
\originalsection{作业9: 2011年11月17日交}
所有 Fourier 变换均采用 $\widehat f(\xi)=\int\ee^{-2\pi\ii x\xi}f(x)\dd x$.
\begin{exercise}\label{ps:9I}
\textbf{原题 I(a--c).} 分别分类 $u_{tt}+u_{tx}+u_{xx}=0$, $u_{tt}+2u_{tx}+u_{xx}=0$, $2u_{tt}-u_{tx}-12u_{xx}=0$.
\end{exercise}
\begin{solution}
主矩阵依次为 $\left(\begin{smallmatrix}1&1/2\\1/2&1\end{smallmatrix}\right)$, $\left(\begin{smallmatrix}1&1\\1&1\end{smallmatrix}\right)$, $\left(\begin{smallmatrix}2&-1/2\\-1/2&-12\end{smallmatrix}\right)$. (a)特征值 $3/2,1/2$ 正, 椭圆型. (b)特征值2、0, 退化抛物型主部; 换 $s=t+x,r=t-x$ 后为 $4u_{ss}=0$, 不同于带时间一阶漂移的热方程. (c)行列式 $-24-1/4<0$, 两特征值异号, 双曲型.
\end{solution}
\begin{exercise}\label{ps:9II}
\textbf{原题 II(a,b).} 定义 $\operatorname{sinc}x=\sin(\pi x)/(\pi x)$, 在0取1. 证光滑, 并求 $\mathbf1_{[-a,a]}$ 的 Fourier 变换, $a>0$.
\end{exercise}
\begin{solution}
(a) $\operatorname{sinc}x=\sum_{k\ge0}(-1)^k(\pi x)^{2k}/(2k+1)!$, 幂级数在紧集可任意逐项求导, 故全线 $C^\infty$.
(b) $\xi\ne0$ 时直接积分得 $\int_{-a}^a\ee^{-2\pi\ii x\xi}\dd x=\sin(2\pi a\xi)/(\pi\xi)=2a\operatorname{sinc}(2a\xi)$. $\xi=0$ 时两边均为2a.
\end{solution}
\begin{exercise}\label{ps:9III}
\textbf{原题 III.} 令 $C_0(\R)=\{u\in C(\R;\C):u(x)\to0\ (x\to\pm\infty)\}$, 范数为 $\sup_x|u(x)|$. 若 $u_j\in C_0(\R)$ 且一致收敛至 $u$, 证明 $u\in C_0$.
\end{exercise}
\begin{solution}
一致极限连续. 给 $\eps>0$, 选 $j$ 使 $\|u-u_j\|_\infty<\eps/2$, 再选 $R$ 使 $|x|>R$ 时 $|u_j(x)|<\eps/2$. 则 $|u(x)|<\eps$, 故两端趋0.
\end{solution}
\begin{exercise}\label{ps:9IV}
\textbf{原题 IV.} $f$ 连续且紧支撑, $\tau_yf(x)=f(x-y)$, 证明 $\|\tau_yf-f\|_1\to0$ 当 $y\to0$.
\end{exercise}
\begin{solution}
$f$ 全线一致连续. 若支撑在 $[-R,R]$, $|y|<1$ 时差支撑于 $[-R-1,R+1]$. 所以范数不超过 $(2R+2)\sup_x|f(x-y)-f(x)|\to0$.
\end{solution}
\begin{exercise}\label{ps:9V}
\textbf{原题 V.} 对同样的 $f$, 证明 $f_t=\Gamma(t)*f\to f$ 于 $L^1$.
\end{exercise}
\begin{solution}
核质量1与 Tonelli 给 $\|f_t-f\|_1\le\int\Gamma(t,y)\|\tau_yf-f\|_1\dd y$. 选 $r>0$ 使 $|y|<r$ 时内范数 $<\eps$, 小位移部分不超过 $\eps$. 其余内范数不超过 $2\|f\|_1$, 核尾 $\int_{|y|\ge r}\Gamma(t,y)\dd y=\int_{|z|\ge r/\sqrt t}\Gamma(1,z)\dd z\to0$. 先定r再令t趋0, 得结论.
\end{solution}
\begin{exercise}\label{ps:9VI}
\textbf{原题 VI(a--e).} 沿热正则化证明上述 $f$ 的 Riemann--Lebesgue 引理: (a)求 $\widehat f_t$; (b)连续有界与误差界; (c)$\widehat f_t\in C_0$; (d)一致趋 $\widehat f$; (e)$\widehat f$ 两端趋0.
\end{exercise}
\begin{solution}
(a)卷积定理及 Gaussian 变换给 $\widehat f_t=\ee^{-4\pi^2t\xi^2}\widehat f$.
(b)对任意 $h\in L^1$, $|\widehat h|\le\|h\|_1$. 对 $\xi_j\to\xi$, 指数逐点收敛并受 $|h|$ 控制, 故变换连续. 特别 $\|\widehat f_t-\widehat f\|_\infty\le\|f_t-f\|_1$.
(c) $\widehat f$ 有界, Gaussian 因子趋0, 故 $\widehat f_t\in C_0$.
(d)由(b)与原题 V, 误差趋0.
(e)用原题 III 的一致闭合性得 $\widehat f\in C_0$. 此题的源假设为连续紧支撑, 一般 $L^1$ 版可通过该类在 $L^1$ 的稠密性和同样的一致误差估计推广.
\end{solution}
\originalsection{作业10: 2011年12月1日交}
\begin{exercise}\label{ps:10I}
\textbf{原题 I(a--d).} $f\in C_c^\infty(\R^n)$, 沿 Fourier 变换重推热核及全空间热初值解. 为排除无穷远异常解, 采用 Schwartz 或连续 $L^2$ 能量解类.
\end{exercise}
\begin{solution}
(a) $\widehat{\Delta u}=-4\pi^2|\xi|^2\widehat u$, 所以 $\partial_t\widehat u=-4\pi^2|\xi|^2\widehat u$, 初值 $\widehat f$.
(b)标量 ODE 得 $\widehat u=\ee^{-4\pi^2t|\xi|^2}\widehat f$.
(c)一维 Gaussian 积分 $\int\ee^{-a\xi^2+2\pi\ii x\xi}\dd\xi=\sqrt{\pi/a}\ee^{-\pi^2x^2/a}$, 取 $a=4\pi^2t$; 各坐标相乘得逆核 $(4\pi t)^{-n/2}\ee^{-|x|^2/(4t)}=\Gamma(t,x)$.
(d)卷积定理和反演给 $u=\Gamma(t)*f$. 公式正时间 Schwartz, 初时以近似恒等收敛. 任意仅光滑且无增长约束的全空间解不能无条件使用此 Fourier 推导.
\end{solution}
\begin{exercise}\label{ps:10II}
\textbf{原题 II(a--c).} 实值 $f\in C_c^\infty(\R)$, 证明 $\int f^2=-2\int xff'$, 推出 $\|f\|_2^2\le2\|xf\|_2\|f'\|_2$, 再得 Heisenberg 不等式 $\|f\|_2^2\le4\pi\|xf\|_2\|\xi\widehat f\|_2$.
\end{exercise}
\begin{solution}
(a)积分 $\partial_x(xf^2)=f^2+2xff'$, 紧支撑给边界项0.
(b)上式取绝对值后用 Cauchy--Schwarz.
(c)导数变换 $\widehat{f'}=2\pi\ii\xi\widehat f$, Plancherel 给 $\|f'\|_2=2\pi\|\xi\widehat f\|_2$, 代入(b). 对非零f两种局部化尺度乘积不能任意小.
\end{solution}
\begin{exercise}\label{ps:10III}
\textbf{原题 III(a--c).} $f\in C_c^\infty(\R;\C)$, 证明 $f(x)\ee^{-2\pi\ii\xi x}$ 的指数展开对全体x及有界 $|\xi|<R$ 绝对一致收敛, 可逐项积分, 且非零 $\widehat f$ 不能紧支撑或在非空开区间消失.
\end{exercise}
\begin{solution}
(a)若 $\supp f\subset[-M,M]$, 每项绝对值不超过 $\|f\|_\infty(2\pi RM)^k/k!$, 其和有限; 支撑外每项为0. Weierstrass 判别保证一致绝对收敛, 和为 $f(x)\ee^{-2\pi\ii\xi x}$.
(b)全体项同处有限支撑, 可逐项积分,
$\widehat f(\xi)=\sum_{k\ge0}\xi^k\int f(x)(-2\pi\ii x)^k\dd x/k!$.
系数绝对值不超过 $\|f\|_1(2\pi M)^k/k!$, 所得级数对有界复 $\xi$ 一致绝对收敛, 故是整函数.
(c)解析函数若在开区间为0, 在区间一点所有导数0, Taylor 级数在那里恒0, 再由相交圆盘延拓到全平面; 在任意中心 $\xi_0$ 把f换为 $\ee^{-2\pi\ii\xi_0x}f$, 同一估计给关于 $\xi-\xi_0$ 的无限半径级数, 因此可直接由零点处全部Taylor系数为0得全域为0. 所以 $\widehat f\equiv0$, 反演得 $f=0$. 非零f的变换不能紧支撑, 因紧支撑外已有开区间零值. 必须保留零函数例外.
\end{solution}
\begin{exercise}\label{ps:10IV}
\textbf{原题 IV(a--c), 含符号勘误.} $\ii\psi_t+\Delta\psi/2=F$, $\psi(0)=\phi$, 推积分因子、频率 Duhamel 式及空间核式. 取 $\phi\in C_c^\infty$, $F\in C([0,T];\Sch)$ 且所需半范数在有限时间段一致受控.
\end{exercise}
\begin{solution}
(a)令 $\omega=2\pi^2|\xi|^2$, 方程为 $\partial_t\widehat\psi=-\ii\omega\widehat\psi-\ii\widehat F$, 所以 $\partial_t(\ee^{\ii\omega t}\widehat\psi)=-\ii\ee^{\ii\omega t}\widehat F$.
(b)积分后乘回得到
\[
\widehat\psi(t,\xi)=\ee^{-\ii\omega t}\widehat\phi(\xi)-\ii\int_0^t\ee^{-\ii\omega(t-s)}\widehat F(s,\xi)\dd s.
\]
原卷(b)源项指数误写正号, 这里改为负号.
(c)反演与第十一章已证的振荡核给
$\psi=K(t)*\phi-\ii\int_0^tK(t-s)*F(s)\dd s$,
$K(t,x)=(2\pi\ii t)^{-n/2}\ee^{\ii|x|^2/(2t)}$.
端点 $t=s$ 用酉群强极限 $U(0)=I$ 解释, 不是代入奇异核. 所选 Schwartz 类允许频率求导和换积分. 更弱的 $L^1_tL^2_x$ 源项用群的 Bochner 积分; 原卷(c)的“形式”核反演不能直接当成任意逐时紧支撑源的定理.
\end{solution}

\originalsection{作业11: 2011年12月8日交}
\begin{exercise}\label{ps:11I}
\textbf{原题 I.} 可微可逆矩阵 $N(q)$, 证明 $(N^{-1})'=-N^{-1}N'N^{-1}$.
\end{exercise}
\begin{solution}
微分 $N^{-1}N=I$ 得 $(N^{-1})'N+N^{-1}N'=0$. 右乘 $N^{-1}$ 即结论; 按分量写为
\[
\partial_q(N^{-1})^{\mu\nu}=-(N^{-1})^{\mu\alpha}(\partial_qN_{\alpha\beta})(N^{-1})^{\beta\nu}.
\]
\end{solution}
\begin{exercise}\label{ps:11II}
\textbf{原题 II(a--d).} 用 Frobenius 范数 $\|A\|=(\sum_{ij}|A_{ij}|^2)^{1/2}$, 证明次乘性; 小矩阵的行列式展开; 有限几何级数; 充分小时的 Neumann 逆级数.
\end{exercise}
\begin{solution}
(a)逐项 Cauchy--Schwarz 给 $|(AB)_{ij}|^2\le(\sum_k|A_{ik}|^2)(\sum_k|B_{kj}|^2)$, 对i,j求和得 $\|AB\|^2\le\|A\|^2\|B\|^2$.
(b)Leibniz 行列式公式中常数项为1, 一次项仅对角 $\sum_iA_{ii}=\operatorname{tr}A$, 其余有限多单项式次数至少2. 因 $|A_{ij}|\le\|A\|$, 在 $\|A\|\le1$ 有
$\det(I+A)=1+\operatorname{tr}A+\|A\|^2R(A)$, $R$ 有界. 原卷首项写矩阵I, 应为标量1.
(c)直接展开相消得 $(I-A)\sum_{k=0}^NA^k=I-A^{N+1}$, 因A的幂互相交换, 左右次序均可.
(d)由(b)及 $|\operatorname{tr}A|\le\sqrt{n+1}\|A\|$ 可选 $\gamma<1$ 使 $\det(I-A)\ne0$. 又(a)给 $\|A^k\|\le\gamma^k$ ($k\ge1$), 逆级数绝对收敛于矩阵S; (c)取极限给 $(I-A)S=S(I-A)=I$, 所以 $S=(I-A)^{-1}$. 实际单用 $\|A\|<1$ 的级数已能直接证明可逆.
\end{solution}
\begin{exercise}\label{ps:11III}
\textbf{原题 III(a--d).} $\mathcal L=-m^{\alpha\beta}\phi_\alpha\phi_\beta/2-V(\phi)$, $V$ 光滑实值. 说明坐标协变, 推场方程, 求应力张量并直接证明无散度.
\end{exercise}
\begin{solution}
(a) $\phi$ 为标量, $\phi_\alpha$ 为协向量, 度量逆按双逆指标变换, 收缩 $m^{\alpha\beta}\phi_\alpha\phi_\beta$ 为标量, $V(\phi)$ 亦然. 任意坐标下度量分量须同步变换; 不能变坐标却仍强行令矩阵diag不变. 不变作用量应配体积密度, 在 Lorentz 坐标中Jacobian绝对值1.
(b)紧支撑变分h给 $\delta S=\int[-\partial^\alpha\phi\partial_\alpha h-V'(\phi)h]=\int(\Box\phi-V'(\phi))h$. 故 $\Box\phi=V'(\phi)$, 即 $\phi_{tt}-\Delta\phi+V'(\phi)=0$.
(c)度量变分或第十二章张量公式给
$T^{\mu\nu}=\partial^\mu\phi\partial^\nu\phi-m^{\mu\nu}(Q/2+V(\phi))$.
(d)乘积求导后二阶交叉项与 $\partial^\nu Q/2$ 抵消,
$\partial_\mu T^{\mu\nu}=(\Box\phi-V'(\phi))\partial^\nu\phi=0$.
若 $V=\kappa^2\phi^2/2$, 场方程为质量参数 $\kappa$ 的 Klein--Gordon 方程.
\end{solution}
\begin{exercise}\label{ps:11IV}
\textbf{原题 IV(a--d), 含勘误.} 设 $\phi$ 为原题 III 场方程的 $C^2$ 解, $T$ 为其应力张量. $Y$ 的流 $\widetilde x$ 保持 Minkowski 度量, $M=\partial\widetilde x/\partial x$. (a)微分 $Mm^{-1}M^T=m^{-1}$ 得 Killing 方程; (b)证 $\partial_\mu(T^{\mu\nu}Y_\nu)=0$ 及其区域积分式; (c)解释守恒量; (d)取 $Y=(-1,0,\ldots)$ 得带势能量守恒.
\end{exercise}
\begin{solution}
(a)在流参数0处 $M=I$, $\partial_\eps M^\mu_\alpha=\partial_\alpha Y^\mu$. 微分后
$m^{\alpha\nu}\partial_\alpha Y^\mu+m^{\mu\alpha}\partial_\alpha Y^\nu=0$.
降指标得 $\partial_\mu Y_\nu+\partial_\nu Y_\mu=0$.
(b)使用上题无散度及对称性,
$\partial_\mu(T^{\mu\nu}Y_\nu)=T^{\mu\nu}\partial_\mu Y_\nu=\tfrac12T^{\mu\nu}(\partial_\mu Y_\nu+\partial_\nu Y_\mu)=0$.
对任何积分存在的有界光滑区域K积分也为0. 原卷用“任意光滑Z”的变分积分恒等式, 实际须Z内部紧支撑或有适当边界条件; 这里直接证目标, 不沿用错误的无边界限制变分.
(c)对 $[t_1,t_2]\times B_R$ 积分散度, 得两时间片积分差等于负侧面通量. 若侧面通量为0或 $R\to\infty$ 时趋0, 则 $\int T^{0\nu}Y_\nu$ 守恒. 这是等距对称产生守恒量的具体形式.
(d) $Y_0=1$, 故流为 $T^{\mu0}$, 时间分量 $T^{00}=(\phi_t^2+|\nabla\phi|^2)/2+V(\phi)$. 取初值紧支撑、有限时间段支撑受控的解并假设 $V(0)=V'(0)=0$, 得
\[
\int\left[\tfrac12(\phi_t^2+|\nabla\phi|^2)+V(\phi)\right](t,x)\dd x
=\int\left[\tfrac12(\phi_t^2+|\nabla\phi|^2)+V(\phi)\right](0,x)\dd x.
\]
原卷右边梯度误留t, 此处全部在0评价; 场方程应回指原(0.0.8), 而非 Lagrangian(0.0.7). 若 $V(0)\ne0$, 必须用 $V(\phi)-V(0)$ 消除无限常数势能; 无穷远条件亦不可省.
\end{solution}

\originalchapter{2011年期中考试及官方解答重构}
2011年10月27日星期四, 五题各20分, 总100分. 原卷要求逻辑完整, 指明所用定理及适用条件. 本章答案以 Jared Speck 官方 \textit{Midterm Exam Solutions} 为依据翻译重构; 纠正一般反向热级数的存在性、零能量除法等条件问题, 不把这些修订说成官方勘误.
\begin{exercise}\label{exam:mI}
\textbf{原题 I(a:8分,b:12分).} 区间 $[0,1]$ 的反向热方程 $u_t=-u_{xx}$, 零 Dirichlet 边值. (a)找全部分离变量解 $u=v(t)w(x)$. (b)光滑初值f两端为0且 $\|f\|_\infty\le\eps$, 讨论 $t=1$ 大小, 与零初值及正向热流比较, 说明适定性.
\end{exercise}
\begin{solution}
(a)零解总成立. 对非零乘积, 分离得 $v'/v=-w''/w=\lambda$. 两端零的非零空间解必须 $\lambda=m^2\pi^2$, $m=1,2,\ldots$, 因 $\lambda=0$ 的仿射解及 $\lambda<0$ 的双曲函数解都被两端零强迫为零. 因此全部非零分离解是 $A\ee^{m^2\pi^2t}\sin(m\pi x)$, A任意非零常数. 官方列 $m\ge0$, 其中 $m=0$ 只是零函数.

(b)取 $f=\eps\sin(m\pi x)$, 得 $\|u(1)\|_\infty=\eps\ee^{m^2\pi^2}$ 可任意大. 为严格给连续依赖失败, 取 $\eps_m=\ee^{-m^2\pi^2/2}$, 则初值范数趋0, $t=1$ 范数趋无穷; 零初值对应零解. 正向热流则模态为 $\ee^{-m^2\pi^2t}$, 最大原理给 $\|u(t)\|_\infty\le\|f\|_\infty$, 连续依赖成立.

官方把一般光滑f的展开 $\sum b_m\ee^{m^2\pi^2t}\sin(m\pi x)$ 直接称解, 此处须修订: 它一般只是形式级数. 若要求在正时间 $\tau$ 属于 $L^2$, 必须 $\sum|b_m|^2\ee^{2m^2\pi^2\tau}<\infty$, 即使系数超多项式衰减也未必满足该 Gaussian 加权条件; 仅端点为零的光滑初值甚至可能只有多项式衰减. 例如 $f=x(1-x)$ 的系数为 $4[1-(-1)^m]/(m^3\pi^3)$, 代入上述加权和立即发散, 所以不存在正时间 $L^2$ 解. 反向问题既可无解, 也在已有解的单频序列上失去连续依赖, 故在通常的初值范数中不适定.
\end{solution}
\begin{exercise}\label{exam:mII}
\textbf{原题 II:20分.} $u\in C^\infty(\R^3)$ 调和且 $|u(x)|\le\sqrt{|x|}$, 证明 $u=0$.
\end{exercise}
\begin{solution}
$u(0)=0$. 在 $B_R$ 置 $v=u+\sqrt R\ge0$, $v(0)=\sqrt R$. 官方使用三维球 Harnack:
\[
\left[\frac{R(R-r)}{(R+r)^2}-1\right]\sqrt R\le u(x)\le
\left[\frac{R(R+r)}{(R-r)^2}-1\right]\sqrt R,
\quad r=|x|<R.
\]
固定r时两系数括号为 $O(r/R)$, 乘 $\sqrt R$ 趋0, 夹逼得 $u(x)=0$.
\end{solution}
\begin{exercise}\label{exam:mIII}
\textbf{原题 III:20分.} $u_t-u_{xx}=-u$ 于 $[0,2]\times[0,1]$, 给定非正初值与两个非正边值, $u\in C^{1,2}$, 证 $u\le0$.
\end{exercise}
\begin{solution}
在紧矩形取最大点. 若在初片或空间端点, 数据给非正. 否则假设最大值正; 空间内点有 $u_{xx}\le0$. 若 $0<t_0<2$ 则 $u_t=0$, 若 $t_0=2$ 则左时间导数 $u_t\ge0$. 两种情形均使 $u_t-u_{xx}\ge0$, 而方程右端 $-u<0$, 矛盾. 这是官方最大点证明, 也可对 $\ee^tu$ 用普通热最大原理.
\end{solution}
\begin{exercise}\label{exam:mIV}
\textbf{原题 IV(a:6分,b:2分,c:12分).} $u\in C^{1,2}([0,\infty)\times[0,1])$, $u_t=u_{xx}$ 于该域, $u_x(t,0)=u_x(t,1)=0$, $u(0)=f$. (a)证 $T(t)=\int_0^1u$ 不变; (b)猜长期状态; (c)严格证明在一个合适范数中收敛.
\end{exercise}
\begin{solution}
(a) $T'=\int u_{xx}=u_x(1)-u_x(0)=0$.
(b)绝热端点无热量流失, 稳态应为常数 $C=T(0)=\int_0^1f$.
(c)按官方能量论证置 $w=u-C$, 则 $\int w=0$, 满足齐次热方程及 Neumann 边值. 连续性和零均值给某点 $x_0$ 有 $w(x_0)=0$, 故 $|w(x)|\le\int_0^1|w_x|\le\|w_x\|_2$, 进而 $\|w\|_2^2\le\|w_x\|_2^2$. 分部积分给 $\partial_t\|w\|_2^2=-2\|w_x\|_2^2\le-2\|w\|_2^2$. 积分得 $\|u(t)-C\|_2\le\ee^{-t}\|f-C\|_2\to0$. 官方后面的“积分(21)”误引空间估计, 应积分其(22)的时间微分不等式; 此处明确使用后者.
\end{solution}
\begin{exercise}\label{exam:mV}
\textbf{原题 V(a:8分,b:12分).} 原条件为 $F\in C^2([0,\infty)\times\R)$, $f\in C^2\cap L^2$, $g\in C^1\cap L^2$, $u\in C^2$, 每个时刻u紧支撑. $-u_{tt}+u_{xx}=F$, 初值f,g, $\|F(t)\|_2\le(1+t^2)^{-1}$. 编者为严格积分补充有限能量及无穷远通量控制, 可采用有限时间段统一紧支撑的光滑解类. 定义 $E^2=\int(u_t^2+u_x^2)$. (a)证 $(E^2)'=-2\int Fu_t$; (b)证 $E(t)\le E(0)+C$.
\end{exercise}
\begin{solution}
(a) 求导并以 $u_{tt}=u_{xx}-F$ 代换, 得 $2\int u_tu_{xx}+u_xu_{xt}-Fu_t$. 前两项分部积分相消, 余 $-2\int Fu_t$. 原卷假设逐时紧支撑; 严格积分论证可取有限时间段统一紧支撑, 或用有限能量截断传递.
(b)Cauchy--Schwarz 给 $(E^2)'\le2E\|F\|_2$. 不能在 $E=0$ 时直接除E. 置 $E_\delta=\sqrt{E^2+\delta^2}$, 得 $E_\delta'\le(E/E_\delta)\|F\|_2\le(1+t^2)^{-1}$. 积分再令 $\delta\downarrow0$, 得 $E(t)\le E(0)+\arctan t\le E(0)+\pi/2$. 这补足官方直接除E的零值情形.
\end{solution}

\originalchapter{2011年期末考试及独立解答}
2011年12月19日星期一. 原卷七题总170分, 要求列明定理和条件. 期末卷实际允许一张8.5\,$\times$\,11英寸纸的正反两面手写笔记, 禁止其他书籍、笔记、计算器; 这与教学大纲的闭卷无笔记说明有差异, 本册以实际试卷说明为准. 官方未公开期末答案; 本章全部解答为 AI 独立编写并数学复核.
\begin{exercise}\label{exam:fI}
\textbf{原题 I:20分.} 一维波方程光滑初值f,g在 $|x|\ge R$ 消失, 证明 $t>0$, $|x|\ge R+t$ 时 $u(t,x)=0$.
\end{exercise}
\begin{solution}
$d'Alembert$ 式 $u=[f(x+t)+f(x-t)]/2+\tfrac12\int_{x-t}^{x+t}g(y)\dd y$. 若 $x\ge R+t$, 则两个端点及整个积分区间都在 $[R,\infty)$; 若 $x\le-R-t$, 则都在 $(-\infty,-R]$. 所有项为0. 等号位置也因数据在 $|x|\ge R$ 为0成立.
\end{solution}
\begin{exercise}\label{exam:fII}
\textbf{原题 II(a,b各10分).} $f=\mathbf1_{[-2,2]}$, 求 Fourier 变换及其 $L^2$ 范数.
\end{exercise}
\begin{solution}
(a)直接积分得 $\widehat f(\xi)=\sin(4\pi\xi)/(\pi\xi)=4\operatorname{sinc}(4\xi)$, 0处取连续值4.
(b)Plancherel 给 $\|\widehat f\|_2=\|f\|_2=\sqrt4=2$. 不需额外计算振荡平方积分.
\end{solution}
\begin{exercise}\label{exam:fIII}
\textbf{原题 III(a:5分,b:10分,c:5分,d:10分).} 全空间波方程初位移 $f\in C_c^\infty(\R^n)$, 初速度0. (a)推频率 ODE; (b)求解; (c)求时间导数与空间梯度变换; (d)只用 Fourier 方法, 不用分部积分, 证 $\|\nabla_{t,x}u(t)\|_2=\|\nabla f\|_2$. 原定义(4)被积函数误写f, 应为 $u(t,x)$.
\end{exercise}
\begin{solution}
(a)空间变换给 $\widehat u_{tt}=-4\pi^2|\xi|^2\widehat u$, $\widehat u(0)=\widehat f$, $\widehat u_t(0)=0$.
(b) $\widehat u=\cos(2\pi t|\xi|)\widehat f$, 在 $\xi=0$ 为常数 $\widehat f(0)$.
(c) $\widehat u_t=-2\pi|\xi|\sin(2\pi t|\xi|)\widehat f$, $\widehat{\nabla u}=2\pi\ii\xi\cos(2\pi t|\xi|)\widehat f$.
(d) 对各分量用 Plancherel, 两平方和为
\[
\int4\pi^2|\xi|^2[\sin^2(2\pi t|\xi|)+\cos^2(2\pi t|\xi|)]|\widehat f|^2\dd\xi
=\|\nabla f\|_2^2.
\]
开平方即得, 全部计算在频率端.
\end{solution}
\begin{exercise}\label{exam:fIV}
\textbf{原题 IV:20分.} 给定 $u_t-u_{xx}=-(t^2+x^2)$ 在 $[0,2]\times[0,1]$ 的光滑解, 初值f. 证明整体最大值不可能在 $(0,2)\times(0,1)$ 取得. 原卷未给边值却称唯一 Cauchy 解, 该唯一性无依据; 所证局部最大值命题无需新增边界条件.
\end{exercise}
\begin{solution}
若内部 $(t_0,x_0)$ 为整体最大点, $u_t=0$, $u_{xx}\le0$, 所以 $u_t-u_{xx}\ge0$. 然而 $t_0>0$ 使右边 $-(t_0^2+x_0^2)<0$, 矛盾.
\end{solution}
\begin{exercise}\label{exam:fV}
\textbf{原题 V(a,b各10分).} 自由 Schrödinger 方程 $\ii\psi_t+\Delta\psi/2=0$, 初值 $\phi\in C_c^\infty(\R^n;\C)$. (a)证 $L^2$ 范数守恒; (b)以此证唯一. 原“所有光滑解”应限制到 Schwartz 解或 $C_tL^2$ 能量解类.
\end{exercise}
\begin{solution}
(a)Schwartz 解允许分部积分,
$\partial_t\int|\psi|^2=2\operatorname{Re}\int\overline\psi(\ii\Delta\psi/2)=-\operatorname{Re}(\ii\int|\nabla\psi|^2)=0$.
也可用 $\widehat\psi=\ee^{-2\pi^2\ii t|\xi|^2}\widehat\phi$ 与 Plancherel, 直接推广至 $C_tL^2$ 的分布解类; 唯一性频率截断论证见定理\ref{thm:schroL2}.
(b) 两个同类解之差w仍满足齐次方程, 初值0. 对w用(a)得 $\|w(t)\|_2=0$, 几乎处处相同; 若两解连续则逐点相同. 仅光滑性不能保证积分可积和无穷远边界项消失, 因此不扩大本结论的解类别.
\end{solution}
\begin{exercise}\label{exam:fVI}
\textbf{原题 VI(a:5分,b:15分,c:5分,d:5分).} $\mathcal L=-Q/2-\phi^4/4$. 求 Euler--Lagrange 方程、应力张量及 $T^{00}$ 正性、张量散度, 并说明 $J^\mu=T^{\mu0}$ 给出的守恒量.
\end{exercise}
\begin{solution}
(a) $\Box\phi=\phi^3$, 即 $\phi_{tt}-\Delta\phi+\phi^3=0$.
(b) $T^{\mu\nu}=\partial^\mu\phi\partial^\nu\phi-m^{\mu\nu}(Q/2+\phi^4/4)$, 故 $T^{00}=(\phi_t^2+|\nabla\phi|^2)/2+\phi^4/4\ge0$.
(c)按乘积法则及混合偏导相等, $\partial_\mu T^{\mu\nu}=(\Box\phi-\phi^3)\partial^\nu\phi=0$ 在场方程上.
(d) $J$ 无散度, $J^i=-\phi_i\phi_t$. 若初值紧支撑且考虑存在期间的受控支撑解, 或有限能量且无穷远通量消失的解, 时空柱积分给
$\int[(\phi_t^2+|\nabla\phi|^2)/2+\phi^4/4](t)=\int[(\phi_t^2+|\nabla\phi|^2)/2+\phi^4/4](0)$.
此能量控制每个时刻时空梯度的空间 $L^2$ 范数和场值的空间 $L^4$ 范数, 但本题没有据此证明任意维数任意初值的非线性全局存在.
\end{solution}
\begin{exercise}\label{exam:fVII}
\textbf{原题 VII(a--f各5分).} 举色散方程、非有限传播初值问题; 写三维 Dirichlet Green 方程及边值; 分类 $-u_{tt}+4u_{tx}-u_{xx}=0$; 定义适定性; 举 $\R^2$ 上线性 Cauchy 不适定问题.
\end{exercise}
\begin{solution}
(a)自由 Schrödinger 方程, $\omega(k)=|k|^2/2$, 群速度随频率变化, 并有 $|t|^{-n/2}$ 的 $L^1$ 至 $L^\infty$ 色散界.
(b)全空间热方程 $u_t=\Delta u$, 非负非零紧支撑初值与处处正热核卷积, 任意 $t>0$ 全空间严格正, 所以没有有限传播速度.
(c)沿原课程 $\Delta$ 号约定, 固定x时 $\Delta_yG(x,y)=\delta_x(y)$ 于分布意义, $G(x,y)=0$ 于 $y\in\partial\Omega$. 本书正文若采用 $-\Delta$ 核则右边仍为 $\delta_x$, 核整体变号. 奇点处不能写普通光滑方程.
(d)矩阵 $\left(\begin{smallmatrix}-1&2\\2&-1\end{smallmatrix}\right)$ 特征值1、$-3$, 异号, 双曲型.
(e)指定数据空间、解空间、时间区间及范数后, 对允许数据存在解、唯一且连续依赖数据, 称 Hadamard 适定. “连续依赖”须明确拓扑, 不能仅说有显式公式.
(f)时空 $\R^2$ 的反向热方程 $u_t=-u_{xx}$, 在全线有界周期光滑数据的 $C^0$ 初值范数中不适定: 初值 $f_k=\ee^{-k^2/2}\sin kx\to0$, 解 $u_k=\ee^{k^2(t-1/2)}\sin kx$, 在 $t=1$ 范数趋无穷. 每个成员确为经典全空间解, 所以这是连续依赖的实际反例.
\end{solution}

\originalchapter{来源映射、范围与缺项}
\originalsection{18份课堂讲义的逐文件映射}
原作者均为 Jared Speck, 原课程 MIT 18.152, Fall 2011. 表中页数是实际 PDF 页序, 各含末尾一页 OCW 来源提示, 不是估计正文页数. 所有18份文件哈希均不同, 无整份重复; 合并的是重复概念和论证, 不删除课程主题. 文件名前缀均为 \texttt{mit18\_152f11\_lec\_}. 原件链接及完整 SHA256 见本目录 \texttt{source-manifest.json}, 原件只位于仓库 sources.
\begin{longtable}{p{.13\textwidth}p{.08\textwidth}p{.35\textwidth}p{.32\textwidth}}
\toprule 文件尾号&页数&原讲次及主题&中文正文位置\\\midrule\endhead
01&8&1: PDE与线性性&第1、13章\\
02&8&2: 热方程与分离变量&第2、3章\\
03&4&3: 热唯一性&第2章\\
04&4&4: 热最大原理&第2、4章\\
05&11&5: 热核及全空间&第4章\\
06&7&6: Laplace/Poisson、平均值&第5、6章\\
07&11&7: Poisson 基本解&第6章\\
08&4&8: Green 函数&第7章\\
09&7&9: Poisson 核、Harnack、Liouville&第5、7章\\
10&7&10: 波导论、一维传播&第8章\\
11&4&11: 球面平均&第8章\\
12&7&12: Kirchhoff、Minkowski&第8、9章\\
13\_14&7&13--14: 几何能量、应力张量&第9章及作业/Bonus\\
15&8&15: 二阶方程分类&第10章\\
16\_18&12&16--18: Fourier、反演、Plancherel&第10章\\
19\_20&10&19--20: Schrödinger&第11章\\
21\_23&11&21--23: Lagrange 场论&第12章及作业11\\
24&6&24: 输运与 Burgers&第13章\\\bottomrule
\end{longtable}
\originalsection{课程考核完整性与答案身份}
官方公开11份作业、1份Bonus、期中卷、期中官方答案和期末卷, 共15个PDF、78页. 大量考试空白答题页不算题目. 11份作业有51个罗马编号条目, 其中39道完整公开大题、11条仅商业教材引用、1条阅读任务. 加Bonus1、期中5、期末7, 收入52道完整公开大题、114个叶条目. 无字母分问的大题按1叶计, 内部多个要求仍全部解答. 三个完全重复小问复用已写证明, 不删除其课程位置; 其余111叶各有独立任务.

作业1--11依序对应原题 I--IV; III--V; I、II、V、VI; VI; I--III; I--III; I--V; I、II; I--VI; I--IV; I--IV. 叶条目数依次为4、3、4、1、3、7、7、7、13、13、13. Bonus原卷两个f保留为 $f_1,f_2$, 共10叶. 期中 I--V共9叶, 期末 I--VII共20叶. 每一道题面注明原题号, 可通过目录进入相应作业或考试. Bonus(a,b,$f_1$)分别完全重复作业8-II(a,b,e), 在Bonus明确回指.

只有期中有官方公开答案, 本册基于官方答案重构并明确条件修订. 其余作业、Bonus、期末的答案为 AI 独立编写及交叉数学审读, 不署为MIT官方答案. 官方作业8的错误迹/守恒式由官方Bonus修正题面交叉核验; Bonus(h)点态等式只在积分后成立. 其余发现的原件误号、符号、初值角点和解类限制在对应题处明示, 不暗示官方已发布勘误. 独立审查报告和逐题覆盖记录位于 \texttt{review/}.
\originalsection{商业教材引用的真实缺项}
下表仅登记官方公开作业中的引用, 未获得完整题面, 不猜造问题、不声称已完成解答. 指定书为 Salsa, \textit{Partial Differential Equations in Action: From Modelling to Theory}, Springer, 2010.
\begin{longtable}{p{.2\textwidth}p{.2\textwidth}p{.12\textwidth}p{.35\textwidth}}
\toprule 作业原号&书中题号&书页&公开补充信息\\\midrule\endhead
2-I&2.1&97&无完整题面\\
2-II&2.3&97&仅给 $L=\pi,U=0$\\
3-III&2.13&99&无完整题面\\
3-IV&2.14&99&允许热核表示的提示\\
4-I&3.1&150&无完整题面\\
4-II&3.2&150&无完整题面\\
4-III&3.3&151&补 $u\in C^3(\Omega)\cap C^1(\overline\Omega)$\\
4-IV&3.4&151&无完整题面\\
4-V&3.8&152&无完整题面\\
5-IV&3.14&153&Green 负性与对称性提示, 非题面\\
5-V&3.21&155&Helmholtz 基本解提示, 非题面\\\bottomrule
\end{longtable}
作业1-V要求阅读书的附录A, 只登记阅读任务. 这些商业书引用不受本项目的OCW资源许可自动授权, 源码包不含该书或猜造的书题.
\originalsection{核心范围与编者补充}
本册覆盖全部18份公开课堂讲义的数学主线: 热初边值、能量与最大原理、Gaussian 基本解及增长受控的全空间问题; 调和平均值、Poisson基本解与Green核、Poisson公式、Harnack和Liouville; 一维和三维显式波传播、二维下降法、锥能量与有限传播; 主部分类、Fourier反演与Plancherel; 自由Schrödinger及色散; Lorentz应力张量、变分原理、等距流与守恒; 线性输运和Burgers首次破裂. 核心结论给证明, 公开课程题逐项中文作答.

编者补充包括: 分布初值和基本解的测试函数表述; 热级数的 $L^2$ 与角点解释; Fejér完备性证明; Poisson能量弱解的Riesz构造; Duhamel换积分条件; Fourier频率截断的能量类唯一性; 含势边界能量条件; Burgers弱激波、Rankine--Hugoniot与基本熵Riemann解. 这些帮助连接已有实分析、泛函和ODE基础, 不把它们标作原课额外公开讲义.

未建立的范围: 一般粗糙域椭圆边界正则性与全套Sobolev理论; 一般非线性PDE的局部/全局存在; 任意Burgers数据的熵解完整存在唯一性; 任意势的自伴算子定义域; Maxwell系统、流体力学系统及高维非线性色散理论; 有限差分、有限元与数值收敛. 原课程商业教材全书及11处题面缺项也不在已覆盖范围.
\originalsection{署名、许可与变更}
原课程材料: Jared Speck, MIT OpenCourseWare, 18.152 Introduction to Partial Differential Equations, Fall 2011. 中文整理为翻译、重排、增补证明、AI题解、条件修订和自绘图. 改编采用 CC BY-NC-SA 4.0, 保留署名、非商业和相同许可要求, 不表示原作者或MIT认可. 逐个PDF核对实际页数、首页作者、末页OCW提示、资源页许可及第三方声明: 共33份, 未检出单独第三方版权声明. 商业教材是明确的非公开来源边界. 权利结论依据归档原件及当时的官方条款, 不概括为所有MIT材料均无例外.

官方入口: \href{https://ocw.mit.edu/courses/18-152-introduction-to-partial-differential-equations-fall-2011/pages/lecture-notes/}{课堂讲义}, \href{https://ocw.mit.edu/courses/18-152-introduction-to-partial-differential-equations-fall-2011/pages/assignments/}{作业与Bonus}, \href{https://ocw.mit.edu/courses/18-152-introduction-to-partial-differential-equations-fall-2011/pages/exams/}{考试和期中答案}, \href{https://ocw.mit.edu/pages/privacy-and-terms-of-use/}{OCW许可条款}.

\begin{thebibliography}{9}
\bibitem{speck} Jared Speck. \textit{18.152 Introduction to Partial Differential Equations, Fall 2011}. MIT OpenCourseWare. \href{https://ocw.mit.edu/courses/18-152-introduction-to-partial-differential-equations-fall-2011/pages/lecture-notes/}{18份课堂讲义目录}. 具体讲次对应见末章.
\bibitem{terms} MIT OpenCourseWare. \href{https://ocw.mit.edu/pages/privacy-and-terms-of-use/}{使用条款}; \href{https://creativecommons.org/licenses/by-nc-sa/4.0/}{CC BY-NC-SA 4.0}. 资源权利及实际页数见 source-manifest.json.
\end{thebibliography}
\end{document}
